% Header: Here are all packages used and some additional definitions
%%%%%%%%%%%%%%%%%%%%%%%%%%%%%%%%%%%%%%%%%%%%%%%%%%%%%%%%%%%%%%%%%%%
\documentclass[11pt,a4paper]{article}

\usepackage[margin=2.5cm]{geometry}
\usepackage[onehalfspacing]{setspace}
\usepackage{graphicx} % zum Einbinden von Graphiken
\usepackage{amsmath,amsthm,amssymb} % Mathematik Umgebung
\usepackage{icomma} % Intelligentes Komma, das den richtigen Abstand zwischen Dezimalzahlen als auch in Formeln waehlt.
%\usepackage[ngerman]{babel} % Deutsche Bezeichnungen bei Inhaltsangabe etc.
\usepackage[T1]{fontenc} % andere Schriftsatzkodierung fuer richtige Silbentrennung bei Umlauten
\usepackage[locale=DE,space-before-unit=true,per-mode=symbol]{siunitx} % Bessere Einheiten
\usepackage{placeins} % Definiert den Befehl \FloatBarrier
\usepackage{xcolor}
\usepackage{tabularray}
\usepackage{longtable}
\usepackage{xspace} % Correct spacing after named taxon commands

\graphicspath{{Bilder/}}
\DeclareSIUnit{\dBm}{dBm}
\DeclareSIUnit[per-mode=reciprocal]\WN{\per\centi\meter}

% ------------------------------------------------------------------
% Taxonomic-rank colors
% ------------------------------------------------------------------
\definecolor{TaxPhylum}{HTML}{C62828} % red
\definecolor{TaxClass}{HTML}{EF6C00}  % orange
\definecolor{TaxOrder}{HTML}{F9A825}  % yellow
\definecolor{TaxFamily}{HTML}{2E7D32} % green
\definecolor{TaxGenus}{HTML}{1565C0}  % blue
\definecolor{TaxSpecies}{HTML}{6A1B9A}% purple

% Taxonomic-rank formatting commands
\newcommand{\taxPhylum}[1]{\textcolor{TaxPhylum}{#1}}
\newcommand{\taxClass}[1]{\textcolor{TaxClass}{#1}}
\newcommand{\taxOrder}[1]{\textcolor{TaxOrder}{#1}}
\newcommand{\taxFamily}[1]{\textcolor{TaxFamily}{#1}}
\newcommand{\taxGenus}[1]{\textcolor{TaxGenus}{#1}}
\newcommand{\taxSpecies}[1]{\textcolor{TaxSpecies}{#1}}

% ------------------------------------------------------------------
% General document commands
% ------------------------------------------------------------------
\newcommand{\up}{\ensuremath{\uparrow}}
\newcommand{\down}{\ensuremath{\downarrow}}

% --- Scaffolding helpers (to be filled in paper-by-paper) ---

% Placeholder for content still to be completed from the source paper.
\newcommand{\tbd}{\textcolor{gray}{\emph{[TBD]}}}

% Stacked reference cell: cohort name + linked citation,
% BioProject ID underneath.
% Usage: \refcell{CohortName}{bibkey}{PRJ...}
\newcommand{\refcell}[3]{%
  \shortstack[l]{#1~\cite{#2}\\[2pt]{\scriptsize #3}}%
}

% ------------------------------------------------------------------
% Bibliography
% ------------------------------------------------------------------
\usepackage[
  backend=biber,
  style=numeric,
  sorting=none,
  autocite=superscript,
  maxnames=10
]{biblatex}

\addbibresource{references.bib}

\AtBeginBibliography{%
  \renewbibmacro*{begentry}{\printnames{author}}%
}

\AtEveryBibitem{\clearfield{note}}
\setlength{\bibhang}{1em}

\defbibenvironment{bibliography}
{\list
  {\printtext[labelnumberwidth]{%
     \printfield{labelprefix}%
     \printfield{labelnumber}}}
  {\setlength{\labelwidth}{\labelnumberwidth}%
   \setlength{\leftmargin}{\bibhang}%
   \setlength{\labelsep}{\biblabelsep}%
   \setlength{\itemsep}{\bibitemsep}%
   \setlength{\parsep}{\bibparsep}%
   \addtolength{\leftmargin}{\labelsep}%
   \setlength{\itemindent}{-\bibhang}%
   \renewcommand*{\makelabel}[1]{\hss##1}%
   \raggedright}}
{\endlist}
{\item}

% ------------------------------------------------------------------
% Hyperlinks and invisible taxonomic hover tooltips
% ------------------------------------------------------------------
\usepackage[
  breaklinks=true,
  colorlinks=true,
  linkcolor=blue,
  urlcolor=blue,
  citecolor=blue
]{hyperref}

\usepackage{pdfcomment}

% Named taxon commands and their full paper-aligned classifications.
% Keep taxa-dictionary.tex in the same Overleaf folder as this file.
% taxa-dictionary.tex
% Paper-specific taxon commands for the Artacho et al. allo-HSCT summary.
%
% Required packages/definitions in the main document:
%   \usepackage{xcolor}
%   \usepackage{xspace}
%   \usepackage{hyperref}
%   \usepackage{pdfcomment}
%   \taxOrder, \taxFamily, \taxGenus, and \taxSpecies
%
% Taxonomy convention:
%   The lineage labels below intentionally follow the legacy/paper-aligned
%   terminology used with the paper's SILVA 138-era 16S workflow. Newer
%   taxonomies may instead use names such as Bacillota/Bacteroidota and may
%   split the former Clostridiales into several orders.
%
% Turn all tooltip annotations off by placing \TaxTooltipsfalse AFTER the
% % taxa-dictionary.tex
% Paper-specific taxon commands for the Artacho et al. allo-HSCT summary.
%
% Required packages/definitions in the main document:
%   \usepackage{xcolor}
%   \usepackage{xspace}
%   \usepackage{hyperref}
%   \usepackage{pdfcomment}
%   \taxOrder, \taxFamily, \taxGenus, and \taxSpecies
%
% Taxonomy convention:
%   The lineage labels below intentionally follow the legacy/paper-aligned
%   terminology used with the paper's SILVA 138-era 16S workflow. Newer
%   taxonomies may instead use names such as Bacillota/Bacteroidota and may
%   split the former Clostridiales into several orders.
%
% Turn all tooltip annotations off by placing \TaxTooltipsfalse AFTER the
% % taxa-dictionary.tex
% Paper-specific taxon commands for the Artacho et al. allo-HSCT summary.
%
% Required packages/definitions in the main document:
%   \usepackage{xcolor}
%   \usepackage{xspace}
%   \usepackage{hyperref}
%   \usepackage{pdfcomment}
%   \taxOrder, \taxFamily, \taxGenus, and \taxSpecies
%
% Taxonomy convention:
%   The lineage labels below intentionally follow the legacy/paper-aligned
%   terminology used with the paper's SILVA 138-era 16S workflow. Newer
%   taxonomies may instead use names such as Bacillota/Bacteroidota and may
%   split the former Clostridiales into several orders.
%
% Turn all tooltip annotations off by placing \TaxTooltipsfalse AFTER the
% \input{taxa-dictionary.tex} line in the main document.

% --------------------------------------------------------------------------
% Tooltip engine
% --------------------------------------------------------------------------
\newif\ifTaxTooltips
\TaxTooltipstrue

% Arguments:
%   #1 = formatted text printed in the manuscript
%   #2 = plain-text fallback for PDF bookmarks
%   #3 = tooltip text
\DeclareRobustCommand{\taxonTooltip}[3]{%
  \texorpdfstring{%
    \ifTaxTooltips
      \pdftooltip{#1}{#3}%
    \else
      #1%
    \fi
  }{#2}%
}

% --------------------------------------------------------------------------
% Reusable lineage strings
% --------------------------------------------------------------------------
\newcommand{\taxLineClostridiales}{%
Domain: Bacteria\textCR%
Phylum: Firmicutes\textCR%
Class: Clostridia\textCR%
Order: Clostridiales%
}

\newcommand{\taxLineRuminococcaceae}{%
Domain: Bacteria\textCR%
Phylum: Firmicutes\textCR%
Class: Clostridia\textCR%
Order: Clostridiales\textCR%
Family: Ruminococcaceae%
}

\newcommand{\taxLineBlautia}{%
Domain: Bacteria\textCR%
Phylum: Firmicutes\textCR%
Class: Clostridia\textCR%
Order: Clostridiales\textCR%
Family: Lachnospiraceae\textCR%
Genus: Blautia\textCR%
Species: not specified%
}

\newcommand{\taxLineFaecalibacterium}{%
Domain: Bacteria\textCR%
Phylum: Firmicutes\textCR%
Class: Clostridia\textCR%
Order: Clostridiales\textCR%
Family: Ruminococcaceae\textCR%
Genus: Faecalibacterium\textCR%
Species: not specified%
}

\newcommand{\taxLineOscillibacter}{%
Domain: Bacteria\textCR%
Phylum: Firmicutes\textCR%
Class: Clostridia\textCR%
Order: Clostridiales\textCR%
Family: Ruminococcaceae\textCR%
Genus: Oscillibacter\textCR%
Species: not specified%
}

\newcommand{\taxLineClostridiumIV}{%
Domain: Bacteria\textCR%
Phylum: Firmicutes\textCR%
Class: Clostridia\textCR%
Order: Clostridiales\textCR%
Family: Ruminococcaceae\textCR%
Genus-level 16S cluster: Clostridium IV\textCR%
Species: not resolved%
}

\newcommand{\taxLineFlavonifractor}{%
Domain: Bacteria\textCR%
Phylum: Firmicutes\textCR%
Class: Clostridia\textCR%
Order: Clostridiales\textCR%
Family: Ruminococcaceae\textCR%
Genus: Flavonifractor\textCR%
Species: not specified%
}

\newcommand{\taxLineAnaeromassilibacillusSp}{%
Domain: Bacteria\textCR%
Phylum: Firmicutes\textCR%
Class: Clostridia\textCR%
Order: Clostridiales\textCR%
Family: Ruminococcaceae\textCR%
Genus: Anaeromassilibacillus\textCR%
Species: unresolved (reported as Anaeromassilibacillus sp.)%
}

\newcommand{\taxLineButyricicoccusSp}{%
Domain: Bacteria\textCR%
Phylum: Firmicutes\textCR%
Class: Clostridia\textCR%
Order: Clostridiales\textCR%
Family: Ruminococcaceae\textCR%
Genus: Butyricicoccus\textCR%
Species: unresolved (reported as Butyricicoccus sp.)%
}

\newcommand{\taxLineLachnoanaerobaculumSp}{%
Domain: Bacteria\textCR%
Phylum: Firmicutes\textCR%
Class: Clostridia\textCR%
Order: Clostridiales\textCR%
Family: Lachnospiraceae\textCR%
Genus: Lachnoanaerobaculum\textCR%
Species: unresolved/unclassified%
}

\newcommand{\taxLineEnterococcusFaecium}{%
Domain: Bacteria\textCR%
Phylum: Firmicutes\textCR%
Class: Bacilli\textCR%
Order: Lactobacillales\textCR%
Family: Enterococcaceae\textCR%
Genus: Enterococcus\textCR%
Species: Enterococcus faecium%
}

\newcommand{\taxLineParabacteroides}{%
Domain: Bacteria\textCR%
Phylum: Bacteroidetes\textCR%
Class: Bacteroidia\textCR%
Order: Bacteroidales\textCR%
Family: Tannerellaceae\textCR%
Genus: Parabacteroides\textCR%
Species: not specified%
}

\newcommand{\taxLineStaphylococcus}{%
Domain: Bacteria\textCR%
Phylum: Firmicutes\textCR%
Class: Bacilli\textCR%
Order: Bacillales\textCR%
Family: Staphylococcaceae\textCR%
Genus: Staphylococcus\textCR%
Species: not specified%
}

\newcommand{\taxLineStaphylococcusSpecies}{%
Domain: Bacteria\textCR%
Phylum: Firmicutes\textCR%
Class: Bacilli\textCR%
Order: Bacillales\textCR%
Family: Staphylococcaceae\textCR%
Genus: Staphylococcus\textCR%
Species discussed: S. aureus; S. epidermidis; S. simulans%
}

\newcommand{\taxLineStaphylococcusAureus}{%
Domain: Bacteria\textCR%
Phylum: Firmicutes\textCR%
Class: Bacilli\textCR%
Order: Bacillales\textCR%
Family: Staphylococcaceae\textCR%
Genus: Staphylococcus\textCR%
Species: Staphylococcus aureus%
}

\newcommand{\taxLineStaphylococcusEpidermidis}{%
Domain: Bacteria\textCR%
Phylum: Firmicutes\textCR%
Class: Bacilli\textCR%
Order: Bacillales\textCR%
Family: Staphylococcaceae\textCR%
Genus: Staphylococcus\textCR%
Species: Staphylococcus epidermidis%
}

\newcommand{\taxLineStaphylococcusSimulans}{%
Domain: Bacteria\textCR%
Phylum: Firmicutes\textCR%
Class: Bacilli\textCR%
Order: Bacillales\textCR%
Family: Staphylococcaceae\textCR%
Genus: Staphylococcus\textCR%
Species: Staphylococcus simulans%
}

\newcommand{\taxLineClostridiumXIVb}{%
Domain: Bacteria\textCR%
Phylum: Firmicutes\textCR%
Class: Clostridia\textCR%
Order: Clostridiales\textCR%
Family: not resolved by this paper-level label\textCR%
Genus-level 16S cluster: Clostridium XIVb\textCR%
Species: not resolved%
}

\newcommand{\taxLineBacteroides}{%
Domain: Bacteria\textCR%
Phylum: Bacteroidetes\textCR%
Class: Bacteroidia\textCR%
Order: Bacteroidales\textCR%
Family: Bacteroidaceae\textCR%
Genus: Bacteroides\textCR%
Species: not specified%
}

\newcommand{\taxLineBacteroidesSpecies}{%
Domain: Bacteria\textCR%
Phylum: Bacteroidetes\textCR%
Class: Bacteroidia\textCR%
Order: Bacteroidales\textCR%
Family: Bacteroidaceae\textCR%
Genus: Bacteroides\textCR%
Species: unspecified/multiple species%
}

\newcommand{\taxLineVeillonellaceae}{%
Domain: Bacteria\textCR%
Phylum: Firmicutes\textCR%
Class: Negativicutes\textCR%
Order: Selenomonadales\textCR%
Family: Veillonellaceae%
}

\newcommand{\taxLineSelenomonadales}{%
Domain: Bacteria\textCR%
Phylum: Firmicutes\textCR%
Class: Negativicutes\textCR%
Order: Selenomonadales%
}

\newcommand{\taxLineParabacteroidesDistasonis}{%
Domain: Bacteria\textCR%
Phylum: Bacteroidetes\textCR%
Class: Bacteroidia\textCR%
Order: Bacteroidales\textCR%
Family: Porphyromonadaceae\textCR%
Genus: Parabacteroides\textCR%
Species: Parabacteroides distasonis%
}

\newcommand{\taxLineBarnesiellaceae}{%
Domain: Bacteria\textCR%
Phylum: Bacteroidetes\textCR%
Class: Bacteroidia\textCR%
Order: Bacteroidales\textCR%
Family: Barnesiellaceae%
}

\newcommand{\taxLineLachnobacterium}{%
Domain: Bacteria\textCR%
Phylum: Firmicutes\textCR%
Class: Clostridia\textCR%
Order: Clostridiales\textCR%
Family: Lachnospiraceae\textCR%
Genus: Lachnobacterium\textCR%
Species: unresolved (reported as Lachnobacterium unknown)%
}

\newcommand{\taxLineHaemophilus}{%
Domain: Bacteria\textCR%
Phylum: Proteobacteria\textCR%
Class: Gammaproteobacteria\textCR%
Order: Pasteurellales\textCR%
Family: Pasteurellaceae\textCR%
Genus: Haemophilus\textCR%
Species: unresolved (reported as Haemophilus unknown)%
}

\newcommand{\taxLineEnterococcus}{%
Domain: Bacteria\textCR%
Phylum: Firmicutes\textCR%
Class: Bacilli\textCR%
Order: Lactobacillales\textCR%
Family: Enterococcaceae\textCR%
Genus: Enterococcus\textCR%
Species: not specified%
}

\newcommand{\taxLineEnterococcusFaecalis}{%
Domain: Bacteria\textCR%
Phylum: Firmicutes\textCR%
Class: Bacilli\textCR%
Order: Lactobacillales\textCR%
Family: Enterococcaceae\textCR%
Genus: Enterococcus\textCR%
Species: Enterococcus faecalis%
}

\newcommand{\taxLineLactobacillus}{%
Domain: Bacteria\textCR%
Phylum: Firmicutes\textCR%
Class: Bacilli\textCR%
Order: Lactobacillales\textCR%
Family: Lactobacillaceae\textCR%
Genus: Lactobacillus\textCR%
Species: not specified%
}

\newcommand{\taxLineSubdoligranulum}{%
Domain: Bacteria\textCR%
Phylum: Firmicutes\textCR%
Class: Clostridia\textCR%
Order: Clostridiales\textCR%
Family: Ruminococcaceae\textCR%
Genus: Subdoligranulum\textCR%
Species: not specified%
}

\newcommand{\taxLineButyricicoccus}{%
Domain: Bacteria\textCR%
Phylum: Firmicutes\textCR%
Class: Clostridia\textCR%
Order: Clostridiales\textCR%
Family: Ruminococcaceae\textCR%
Genus: Butyricicoccus\textCR%
Species: not specified%
}

\newcommand{\taxLineLachnoclostridium}{%
Domain: Bacteria\textCR%
Phylum: Firmicutes\textCR%
Class: Clostridia\textCR%
Order: Clostridiales\textCR%
Family: Lachnospiraceae\textCR%
Genus: Lachnoclostridium\textCR%
Species: not specified%
}

\newcommand{\taxLineCoprobacter}{%
Domain: Bacteria\textCR%
Phylum: Bacteroidetes\textCR%
Class: Bacteroidia\textCR%
Order: Bacteroidales\textCR%
Family: Porphyromonadaceae\textCR%
Genus: Coprobacter\textCR%
Species: not specified%
}

\newcommand{\taxLineClostridiumInnocuumGroup}{%
Domain: Bacteria\textCR%
Phylum: Firmicutes\textCR%
Class: Erysipelotrichia\textCR%
Order: Erysipelotrichales\textCR%
Family: Erysipelotrichaceae\textCR%
Genus-level cluster: Clostridium (innocuum group)\textCR%
Species: not resolved%
}

\newcommand{\taxLineAlistipes}{%
Domain: Bacteria\textCR%
Phylum: Bacteroidetes\textCR%
Class: Bacteroidia\textCR%
Order: Bacteroidales\textCR%
Family: Rikenellaceae\textCR%
Genus: Alistipes\textCR%
Species: not specified%
}

\newcommand{\taxLineCampylobacter}{%
Domain: Bacteria\textCR%
Phylum: Proteobacteria\textCR%
Class: Epsilonproteobacteria\textCR%
Order: Campylobacterales\textCR%
Family: Campylobacteraceae\textCR%
Genus: Campylobacter\textCR%
Species: not specified%
}

\newcommand{\taxLineLachnospiraceae}{%
Domain: Bacteria\textCR%
Phylum: Firmicutes\textCR%
Class: Clostridia\textCR%
Order: Clostridiales\textCR%
Family: Lachnospiraceae%
}

\newcommand{\taxLineDorea}{%
Domain: Bacteria\textCR%
Phylum: Firmicutes\textCR%
Class: Clostridia\textCR%
Order: Clostridiales\textCR%
Family: Lachnospiraceae\textCR%
Genus: Dorea\textCR%
Species: not specified%
}

\newcommand{\taxLineOscillospira}{%
Domain: Bacteria\textCR%
Phylum: Firmicutes\textCR%
Class: Clostridia\textCR%
Order: Clostridiales\textCR%
Family: Ruminococcaceae\textCR%
Genus: Oscillospira\textCR%
Species: not specified%
}

\newcommand{\taxLineBacteroidaceae}{%
Domain: Bacteria\textCR%
Phylum: Bacteroidetes\textCR%
Class: Bacteroidia\textCR%
Order: Bacteroidales\textCR%
Family: Bacteroidaceae%
}

\newcommand{\taxLineStreptococcus}{%
Domain: Bacteria\textCR%
Phylum: Firmicutes\textCR%
Class: Bacilli\textCR%
Order: Lactobacillales\textCR%
Family: Streptococcaceae\textCR%
Genus: Streptococcus\textCR%
Species: not specified%
}

\newcommand{\taxLineEnterococcaceae}{%
Domain: Bacteria\textCR%
Phylum: Firmicutes\textCR%
Class: Bacilli\textCR%
Order: Lactobacillales\textCR%
Family: Enterococcaceae%
}

\newcommand{\taxLineLactobacillaceae}{%
Domain: Bacteria\textCR%
Phylum: Firmicutes\textCR%
Class: Bacilli\textCR%
Order: Lactobacillales\textCR%
Family: Lactobacillaceae%
}

\newcommand{\taxLineEnterobacteriaceae}{%
Domain: Bacteria\textCR%
Phylum: Proteobacteria\textCR%
Class: Gammaproteobacteria\textCR%
Order: Enterobacteriales\textCR%
Family: Enterobacteriaceae%
}

\newcommand{\taxLineStaphylococcaceae}{%
Domain: Bacteria\textCR%
Phylum: Firmicutes\textCR%
Class: Bacilli\textCR%
Order: Bacillales\textCR%
Family: Staphylococcaceae%
}

% --------------------------------------------------------------------------
% New Lineages added for Post-Transplant Multimodal Summary
% --------------------------------------------------------------------------

\newcommand{\taxLineClostridia}{%
Domain: Bacteria\textCR%
Phylum: Firmicutes\textCR%
Class: Clostridia\textCR%
Order: not specified%
}

\newcommand{\taxLineBacteroidia}{%
Domain: Bacteria\textCR%
Phylum: Bacteroidetes\textCR%
Class: Bacteroidia\textCR%
Order: not specified%
}

\newcommand{\taxLineGammaproteobacteria}{%
Domain: Bacteria\textCR%
Phylum: Proteobacteria\textCR%
Class: Gammaproteobacteria\textCR%
Order: not specified%
}

\newcommand{\taxLineStreptococcusSalivarius}{%
Domain: Bacteria\textCR%
Phylum: Firmicutes\textCR%
Class: Bacilli\textCR%
Order: Lactobacillales\textCR%
Family: Streptococcaceae\textCR%
Genus: Streptococcus\textCR%
Species: Streptococcus salivarius%
}

\newcommand{\taxLineFaecalibacteriumPrausnitzii}{%
Domain: Bacteria\textCR%
Phylum: Firmicutes\textCR%
Class: Clostridia\textCR%
Order: Clostridiales\textCR%
Family: Ruminococcaceae\textCR%
Genus: Faecalibacterium\textCR%
Species: Faecalibacterium prausnitzii%
}

\newcommand{\taxLineBlautiaWexlerae}{%
Domain: Bacteria\textCR%
Phylum: Firmicutes\textCR%
Class: Clostridia\textCR%
Order: Clostridiales\textCR%
Family: Lachnospiraceae\textCR%
Genus: Blautia\textCR%
Species: Blautia wexlerae%
}

\newcommand{\taxLineBacteroidesOvatus}{%
Domain: Bacteria\textCR%
Phylum: Bacteroidetes\textCR%
Class: Bacteroidia\textCR%
Order: Bacteroidales\textCR%
Family: Bacteroidaceae\textCR%
Genus: Bacteroides\textCR%
Species: Bacteroides ovatus%
}

% --------------------------------------------------------------------------
% Named manuscript commands
% --------------------------------------------------------------------------

% Order and family
\DeclareRobustCommand{\taxClostridiales}{%
  \taxonTooltip{\taxOrder{Clostridiales}}{Clostridiales}{\taxLineClostridiales}\xspace}

\DeclareRobustCommand{\taxRuminococcaceae}{%
  \taxonTooltip{\taxFamily{Ruminococcaceae}}{Ruminococcaceae}{\taxLineRuminococcaceae}\xspace}

\DeclareRobustCommand{\taxVeillonellaceae}{%
  \taxonTooltip{\taxFamily{Veillonellaceae}}{Veillonellaceae}{\taxLineVeillonellaceae}\xspace}

\DeclareRobustCommand{\taxSelenomonadales}{%
  \taxonTooltip{\taxOrder{Selenomonadales}}{Selenomonadales}{\taxLineSelenomonadales}\xspace}

% Genera and genus-level 16S clusters
\DeclareRobustCommand{\taxBlautia}{%
  \taxonTooltip{\taxGenus{\textit{Blautia}}}{Blautia}{\taxLineBlautia}\xspace}

\DeclareRobustCommand{\taxFaecalibacterium}{%
  \taxonTooltip{\taxGenus{\textit{Faecalibacterium}}}{Faecalibacterium}{\taxLineFaecalibacterium}\xspace}

\DeclareRobustCommand{\taxOscillibacter}{%
  \taxonTooltip{\taxGenus{\textit{Oscillibacter}}}{Oscillibacter}{\taxLineOscillibacter}\xspace}

\DeclareRobustCommand{\taxClostridiumIV}{%
  \taxonTooltip{\taxGenus{\textit{Clostridium} IV}}{Clostridium IV}{\taxLineClostridiumIV}\xspace}

\DeclareRobustCommand{\taxFlavonifractor}{%
  \taxonTooltip{\taxGenus{\textit{Flavonifractor}}}{Flavonifractor}{\taxLineFlavonifractor}\xspace}

\DeclareRobustCommand{\taxParabacteroides}{%
  \taxonTooltip{\taxGenus{\textit{Parabacteroides}}}{Parabacteroides}{\taxLineParabacteroides}\xspace}

\DeclareRobustCommand{\taxStaphylococcus}{%
  \taxonTooltip{\taxGenus{\textit{Staphylococcus}}}{Staphylococcus}{\taxLineStaphylococcus}\xspace}

\DeclareRobustCommand{\taxClostridiumXIVb}{%
  \taxonTooltip{\taxGenus{\textit{Clostridium} XIVb}}{Clostridium XIVb}{\taxLineClostridiumXIVb}\xspace}

\DeclareRobustCommand{\taxBacteroides}{%
  \taxonTooltip{\taxGenus{\textit{Bacteroides}}}{Bacteroides}{\taxLineBacteroides}\xspace}

% Species-level and unresolved species-level features
\DeclareRobustCommand{\taxAnaeromassilibacillusSp}{%
  \taxonTooltip{\taxSpecies{\textit{Anaeromassilibacillus} sp.}}{Anaeromassilibacillus sp.}{\taxLineAnaeromassilibacillusSp}\xspace}

\DeclareRobustCommand{\taxButyricicoccusSp}{%
  \taxonTooltip{\taxSpecies{\textit{Butyricicoccus} sp.}}{Butyricicoccus sp.}{\taxLineButyricicoccusSp}\xspace}

% Same shotgun feature, displayed without "sp." in the severity bullet.
\DeclareRobustCommand{\taxButyricicoccusBare}{%
  \taxonTooltip{\taxSpecies{\textit{Butyricicoccus}}}{Butyricicoccus}{\taxLineButyricicoccusSp}\xspace}

\DeclareRobustCommand{\taxLachnoanaerobaculumSpecies}{%
  \taxonTooltip{\taxSpecies{unclassified \textit{Lachnoanaerobaculum} species}}{unclassified Lachnoanaerobaculum species}{\taxLineLachnoanaerobaculumSp}\xspace}

\DeclareRobustCommand{\taxEnterococcusFaecium}{%
  \taxonTooltip{\taxSpecies{\textit{Enterococcus faecium}}}{Enterococcus faecium}{\taxLineEnterococcusFaecium}\xspace}

\DeclareRobustCommand{\taxEfaecium}{%
  \taxonTooltip{\taxSpecies{\textit{E. faecium}}}{E. faecium}{\taxLineEnterococcusFaecium}\xspace}

\DeclareRobustCommand{\taxStaphylococcusSpecies}{%
  \taxonTooltip{\taxSpecies{\textit{Staphylococcus} species}}{Staphylococcus species}{\taxLineStaphylococcusSpecies}\xspace}

\DeclareRobustCommand{\taxStaphylococcusSpp}{%
  \taxonTooltip{\taxSpecies{\textit{Staphylococcus} spp.}}{Staphylococcus spp.}{\taxLineStaphylococcusSpecies}\xspace}

\DeclareRobustCommand{\taxStaphylococcusAureus}{%
  \taxonTooltip{\taxSpecies{\textit{Staphylococcus aureus}}}{Staphylococcus aureus}{\taxLineStaphylococcusAureus}\xspace}

\DeclareRobustCommand{\taxSaureus}{%
  \taxonTooltip{\taxSpecies{\textit{S. aureus}}}{S. aureus}{\taxLineStaphylococcusAureus}\xspace}

\DeclareRobustCommand{\taxStaphylococcusEpidermidis}{%
  \taxonTooltip{\taxSpecies{\textit{Staphylococcus epidermidis}}}{Staphylococcus epidermidis}{\taxLineStaphylococcusEpidermidis}\xspace}

\DeclareRobustCommand{\taxSepidermidis}{%
  \taxonTooltip{\taxSpecies{\textit{S. epidermidis}}}{S. epidermidis}{\taxLineStaphylococcusEpidermidis}\xspace}

\DeclareRobustCommand{\taxStaphylococcusSimulans}{%
  \taxonTooltip{\taxSpecies{\textit{Staphylococcus simulans}}}{Staphylococcus simulans}{\taxLineStaphylococcusSimulans}\xspace}

\DeclareRobustCommand{\taxSsimulans}{%
  \taxonTooltip{\taxSpecies{\textit{S. simulans}}}{S. simulans}{\taxLineStaphylococcusSimulans}\xspace}

\DeclareRobustCommand{\taxBacteroidesSpecies}{%
  \taxonTooltip{\taxSpecies{\textit{Bacteroides} species}}{Bacteroides species}{\taxLineBacteroidesSpecies}\xspace}

\DeclareRobustCommand{\taxParabacteroidesDistasonis}{%
  \taxonTooltip{\taxSpecies{\textit{Parabacteroides distasonis}}}{Parabacteroides distasonis}{\taxLineParabacteroidesDistasonis}\xspace}

\DeclareRobustCommand{\taxBarnesiellaceaeSpp}{%
  \taxonTooltip{\taxFamily{Barnesiellaceae spp.}}{Barnesiellaceae spp.}{\taxLineBarnesiellaceae}\xspace}

\DeclareRobustCommand{\taxLachnobacteriumSpp}{%
  \taxonTooltip{\taxSpecies{\textit{Lachnobacterium} spp.}}{Lachnobacterium spp.}{\taxLineLachnobacterium}\xspace}

\DeclareRobustCommand{\taxHaemophilusParainfluenzae}{%
  \taxonTooltip{\taxSpecies{\textit{Haemophilus parainfluenzae}}}{Haemophilus parainfluenzae}{\taxLineHaemophilus}\xspace}

\DeclareRobustCommand{\taxEnterococcus}{%
  \taxonTooltip{\taxGenus{\textit{Enterococcus}}}{Enterococcus}{\taxLineEnterococcus}\xspace}

\DeclareRobustCommand{\taxEnterococcusFaecalis}{%
  \taxonTooltip{\taxSpecies{\textit{Enterococcus faecalis}}}{Enterococcus faecalis}{\taxLineEnterococcusFaecalis}\xspace}

\DeclareRobustCommand{\taxEfaecalis}{%
  \taxonTooltip{\taxSpecies{\textit{E. faecalis}}}{E. faecalis}{\taxLineEnterococcusFaecalis}\xspace}

\DeclareRobustCommand{\taxLactobacillus}{%
  \taxonTooltip{\taxGenus{\textit{Lactobacillus}}}{Lactobacillus}{\taxLineLactobacillus}\xspace}

\DeclareRobustCommand{\taxSubdoligranulum}{%
  \taxonTooltip{\taxGenus{\textit{Subdoligranulum}}}{Subdoligranulum}{\taxLineSubdoligranulum}\xspace}

\DeclareRobustCommand{\taxButyricicoccus}{%
  \taxonTooltip{\taxGenus{\textit{Butyricicoccus}}}{Butyricicoccus}{\taxLineButyricicoccus}\xspace}

\DeclareRobustCommand{\taxLachnoclostridium}{%
  \taxonTooltip{\taxGenus{\textit{Lachnoclostridium}}}{Lachnoclostridium}{\taxLineLachnoclostridium}\xspace}

\DeclareRobustCommand{\taxCoprobacter}{%
  \taxonTooltip{\taxGenus{\textit{Coprobacter}}}{Coprobacter}{\taxLineCoprobacter}\xspace}

\DeclareRobustCommand{\taxClostridiumInnocuumGroup}{%
  \taxonTooltip{\taxGenus{\textit{Clostridium} (innocuum group)}}{Clostridium (innocuum group)}{\taxLineClostridiumInnocuumGroup}\xspace}

\DeclareRobustCommand{\taxAlistipes}{%
  \taxonTooltip{\taxGenus{\textit{Alistipes}}}{Alistipes}{\taxLineAlistipes}\xspace}

\DeclareRobustCommand{\taxCampylobacter}{%
  \taxonTooltip{\taxGenus{\textit{Campylobacter}}}{Campylobacter}{\taxLineCampylobacter}\xspace}

\DeclareRobustCommand{\taxLachnospiraceae}{%
  \taxonTooltip{\taxFamily{Lachnospiraceae}}{Lachnospiraceae}{\taxLineLachnospiraceae}\xspace}

\DeclareRobustCommand{\taxDorea}{%
  \taxonTooltip{\taxGenus{\textit{Dorea}}}{Dorea}{\taxLineDorea}\xspace}

\DeclareRobustCommand{\taxOscillospira}{%
  \taxonTooltip{\taxGenus{\textit{Oscillospira}}}{Oscillospira}{\taxLineOscillospira}\xspace}

\DeclareRobustCommand{\taxBacteroidaceae}{%
  \taxonTooltip{\taxFamily{Bacteroidaceae}}{Bacteroidaceae}{\taxLineBacteroidaceae}\xspace}

\DeclareRobustCommand{\taxStreptococcus}{%
  \taxonTooltip{\taxGenus{\textit{Streptococcus}}}{Streptococcus}{\taxLineStreptococcus}\xspace}

\DeclareRobustCommand{\taxEnterococcaceae}{%
  \taxonTooltip{\taxFamily{Enterococcaceae}}{Enterococcaceae}{\taxLineEnterococcaceae}\xspace}

\DeclareRobustCommand{\taxLactobacillaceae}{%
  \taxonTooltip{\taxFamily{Lactobacillaceae}}{Lactobacillaceae}{\taxLineLactobacillaceae}\xspace}

\DeclareRobustCommand{\taxEnterobacteriaceae}{%
  \taxonTooltip{\taxFamily{Enterobacteriaceae}}{Enterobacteriaceae}{\taxLineEnterobacteriaceae}\xspace}

\DeclareRobustCommand{\taxStaphylococcaceae}{%
  \taxonTooltip{\taxFamily{Staphylococcaceae}}{Staphylococcaceae}{\taxLineStaphylococcaceae}\xspace}

\DeclareRobustCommand{\taxEnterococcusSpp}{%
  \taxonTooltip{\taxSpecies{\textit{Enterococcus} spp.}}{Enterococcus spp.}{\taxLineEnterococcus}\xspace}

\DeclareRobustCommand{\taxLactobacillusSpp}{%
  \taxonTooltip{\taxSpecies{\textit{Lactobacillus} spp.}}{Lactobacillus spp.}{\taxLineLactobacillus}\xspace}

% --------------------------------------------------------------------------
% New Commands added for Post-Transplant Multimodal Summary
% --------------------------------------------------------------------------

\DeclareRobustCommand{\taxClostridia}{%
  \taxonTooltip{\taxClass{Clostridia}}{Clostridia}{\taxLineClostridia}\xspace}

\DeclareRobustCommand{\taxBacteroidia}{%
  \taxonTooltip{\taxClass{Bacteroidia}}{Bacteroidia}{\taxLineBacteroidia}\xspace}

\DeclareRobustCommand{\taxGammaproteobacteria}{%
  \taxonTooltip{\taxClass{Gammaproteobacteria}}{Gammaproteobacteria}{\taxLineGammaproteobacteria}\xspace}

\DeclareRobustCommand{\taxStreptococcusSalivarius}{%
  \taxonTooltip{\taxSpecies{\textit{Streptococcus salivarius}}}{Streptococcus salivarius}{\taxLineStreptococcusSalivarius}\xspace}

\DeclareRobustCommand{\taxFaecalibacteriumPrausnitzii}{%
  \taxonTooltip{\taxSpecies{\textit{Faecalibacterium prausnitzii}}}{Faecalibacterium prausnitzii}{\taxLineFaecalibacteriumPrausnitzii}\xspace}

\DeclareRobustCommand{\taxBlautiaWexlerae}{%
  \taxonTooltip{\taxSpecies{\textit{Blautia wexlerae}}}{Blautia wexlerae}{\taxLineBlautiaWexlerae}\xspace}

\DeclareRobustCommand{\taxBacteroidesOvatus}{%
  \taxonTooltip{\taxSpecies{\textit{Bacteroides ovatus}}}{Bacteroides ovatus}{\taxLineBacteroidesOvatus}\xspace} line in the main document.

% --------------------------------------------------------------------------
% Tooltip engine
% --------------------------------------------------------------------------
\newif\ifTaxTooltips
\TaxTooltipstrue

% Arguments:
%   #1 = formatted text printed in the manuscript
%   #2 = plain-text fallback for PDF bookmarks
%   #3 = tooltip text
\DeclareRobustCommand{\taxonTooltip}[3]{%
  \texorpdfstring{%
    \ifTaxTooltips
      \pdftooltip{#1}{#3}%
    \else
      #1%
    \fi
  }{#2}%
}

% --------------------------------------------------------------------------
% Reusable lineage strings
% --------------------------------------------------------------------------
\newcommand{\taxLineClostridiales}{%
Domain: Bacteria\textCR%
Phylum: Firmicutes\textCR%
Class: Clostridia\textCR%
Order: Clostridiales%
}

\newcommand{\taxLineRuminococcaceae}{%
Domain: Bacteria\textCR%
Phylum: Firmicutes\textCR%
Class: Clostridia\textCR%
Order: Clostridiales\textCR%
Family: Ruminococcaceae%
}

\newcommand{\taxLineBlautia}{%
Domain: Bacteria\textCR%
Phylum: Firmicutes\textCR%
Class: Clostridia\textCR%
Order: Clostridiales\textCR%
Family: Lachnospiraceae\textCR%
Genus: Blautia\textCR%
Species: not specified%
}

\newcommand{\taxLineFaecalibacterium}{%
Domain: Bacteria\textCR%
Phylum: Firmicutes\textCR%
Class: Clostridia\textCR%
Order: Clostridiales\textCR%
Family: Ruminococcaceae\textCR%
Genus: Faecalibacterium\textCR%
Species: not specified%
}

\newcommand{\taxLineOscillibacter}{%
Domain: Bacteria\textCR%
Phylum: Firmicutes\textCR%
Class: Clostridia\textCR%
Order: Clostridiales\textCR%
Family: Ruminococcaceae\textCR%
Genus: Oscillibacter\textCR%
Species: not specified%
}

\newcommand{\taxLineClostridiumIV}{%
Domain: Bacteria\textCR%
Phylum: Firmicutes\textCR%
Class: Clostridia\textCR%
Order: Clostridiales\textCR%
Family: Ruminococcaceae\textCR%
Genus-level 16S cluster: Clostridium IV\textCR%
Species: not resolved%
}

\newcommand{\taxLineFlavonifractor}{%
Domain: Bacteria\textCR%
Phylum: Firmicutes\textCR%
Class: Clostridia\textCR%
Order: Clostridiales\textCR%
Family: Ruminococcaceae\textCR%
Genus: Flavonifractor\textCR%
Species: not specified%
}

\newcommand{\taxLineAnaeromassilibacillusSp}{%
Domain: Bacteria\textCR%
Phylum: Firmicutes\textCR%
Class: Clostridia\textCR%
Order: Clostridiales\textCR%
Family: Ruminococcaceae\textCR%
Genus: Anaeromassilibacillus\textCR%
Species: unresolved (reported as Anaeromassilibacillus sp.)%
}

\newcommand{\taxLineButyricicoccusSp}{%
Domain: Bacteria\textCR%
Phylum: Firmicutes\textCR%
Class: Clostridia\textCR%
Order: Clostridiales\textCR%
Family: Ruminococcaceae\textCR%
Genus: Butyricicoccus\textCR%
Species: unresolved (reported as Butyricicoccus sp.)%
}

\newcommand{\taxLineLachnoanaerobaculumSp}{%
Domain: Bacteria\textCR%
Phylum: Firmicutes\textCR%
Class: Clostridia\textCR%
Order: Clostridiales\textCR%
Family: Lachnospiraceae\textCR%
Genus: Lachnoanaerobaculum\textCR%
Species: unresolved/unclassified%
}

\newcommand{\taxLineEnterococcusFaecium}{%
Domain: Bacteria\textCR%
Phylum: Firmicutes\textCR%
Class: Bacilli\textCR%
Order: Lactobacillales\textCR%
Family: Enterococcaceae\textCR%
Genus: Enterococcus\textCR%
Species: Enterococcus faecium%
}

\newcommand{\taxLineParabacteroides}{%
Domain: Bacteria\textCR%
Phylum: Bacteroidetes\textCR%
Class: Bacteroidia\textCR%
Order: Bacteroidales\textCR%
Family: Tannerellaceae\textCR%
Genus: Parabacteroides\textCR%
Species: not specified%
}

\newcommand{\taxLineStaphylococcus}{%
Domain: Bacteria\textCR%
Phylum: Firmicutes\textCR%
Class: Bacilli\textCR%
Order: Bacillales\textCR%
Family: Staphylococcaceae\textCR%
Genus: Staphylococcus\textCR%
Species: not specified%
}

\newcommand{\taxLineStaphylococcusSpecies}{%
Domain: Bacteria\textCR%
Phylum: Firmicutes\textCR%
Class: Bacilli\textCR%
Order: Bacillales\textCR%
Family: Staphylococcaceae\textCR%
Genus: Staphylococcus\textCR%
Species discussed: S. aureus; S. epidermidis; S. simulans%
}

\newcommand{\taxLineStaphylococcusAureus}{%
Domain: Bacteria\textCR%
Phylum: Firmicutes\textCR%
Class: Bacilli\textCR%
Order: Bacillales\textCR%
Family: Staphylococcaceae\textCR%
Genus: Staphylococcus\textCR%
Species: Staphylococcus aureus%
}

\newcommand{\taxLineStaphylococcusEpidermidis}{%
Domain: Bacteria\textCR%
Phylum: Firmicutes\textCR%
Class: Bacilli\textCR%
Order: Bacillales\textCR%
Family: Staphylococcaceae\textCR%
Genus: Staphylococcus\textCR%
Species: Staphylococcus epidermidis%
}

\newcommand{\taxLineStaphylococcusSimulans}{%
Domain: Bacteria\textCR%
Phylum: Firmicutes\textCR%
Class: Bacilli\textCR%
Order: Bacillales\textCR%
Family: Staphylococcaceae\textCR%
Genus: Staphylococcus\textCR%
Species: Staphylococcus simulans%
}

\newcommand{\taxLineClostridiumXIVb}{%
Domain: Bacteria\textCR%
Phylum: Firmicutes\textCR%
Class: Clostridia\textCR%
Order: Clostridiales\textCR%
Family: not resolved by this paper-level label\textCR%
Genus-level 16S cluster: Clostridium XIVb\textCR%
Species: not resolved%
}

\newcommand{\taxLineBacteroides}{%
Domain: Bacteria\textCR%
Phylum: Bacteroidetes\textCR%
Class: Bacteroidia\textCR%
Order: Bacteroidales\textCR%
Family: Bacteroidaceae\textCR%
Genus: Bacteroides\textCR%
Species: not specified%
}

\newcommand{\taxLineBacteroidesSpecies}{%
Domain: Bacteria\textCR%
Phylum: Bacteroidetes\textCR%
Class: Bacteroidia\textCR%
Order: Bacteroidales\textCR%
Family: Bacteroidaceae\textCR%
Genus: Bacteroides\textCR%
Species: unspecified/multiple species%
}

\newcommand{\taxLineVeillonellaceae}{%
Domain: Bacteria\textCR%
Phylum: Firmicutes\textCR%
Class: Negativicutes\textCR%
Order: Selenomonadales\textCR%
Family: Veillonellaceae%
}

\newcommand{\taxLineSelenomonadales}{%
Domain: Bacteria\textCR%
Phylum: Firmicutes\textCR%
Class: Negativicutes\textCR%
Order: Selenomonadales%
}

\newcommand{\taxLineParabacteroidesDistasonis}{%
Domain: Bacteria\textCR%
Phylum: Bacteroidetes\textCR%
Class: Bacteroidia\textCR%
Order: Bacteroidales\textCR%
Family: Porphyromonadaceae\textCR%
Genus: Parabacteroides\textCR%
Species: Parabacteroides distasonis%
}

\newcommand{\taxLineBarnesiellaceae}{%
Domain: Bacteria\textCR%
Phylum: Bacteroidetes\textCR%
Class: Bacteroidia\textCR%
Order: Bacteroidales\textCR%
Family: Barnesiellaceae%
}

\newcommand{\taxLineLachnobacterium}{%
Domain: Bacteria\textCR%
Phylum: Firmicutes\textCR%
Class: Clostridia\textCR%
Order: Clostridiales\textCR%
Family: Lachnospiraceae\textCR%
Genus: Lachnobacterium\textCR%
Species: unresolved (reported as Lachnobacterium unknown)%
}

\newcommand{\taxLineHaemophilus}{%
Domain: Bacteria\textCR%
Phylum: Proteobacteria\textCR%
Class: Gammaproteobacteria\textCR%
Order: Pasteurellales\textCR%
Family: Pasteurellaceae\textCR%
Genus: Haemophilus\textCR%
Species: unresolved (reported as Haemophilus unknown)%
}

\newcommand{\taxLineEnterococcus}{%
Domain: Bacteria\textCR%
Phylum: Firmicutes\textCR%
Class: Bacilli\textCR%
Order: Lactobacillales\textCR%
Family: Enterococcaceae\textCR%
Genus: Enterococcus\textCR%
Species: not specified%
}

\newcommand{\taxLineEnterococcusFaecalis}{%
Domain: Bacteria\textCR%
Phylum: Firmicutes\textCR%
Class: Bacilli\textCR%
Order: Lactobacillales\textCR%
Family: Enterococcaceae\textCR%
Genus: Enterococcus\textCR%
Species: Enterococcus faecalis%
}

\newcommand{\taxLineLactobacillus}{%
Domain: Bacteria\textCR%
Phylum: Firmicutes\textCR%
Class: Bacilli\textCR%
Order: Lactobacillales\textCR%
Family: Lactobacillaceae\textCR%
Genus: Lactobacillus\textCR%
Species: not specified%
}

\newcommand{\taxLineSubdoligranulum}{%
Domain: Bacteria\textCR%
Phylum: Firmicutes\textCR%
Class: Clostridia\textCR%
Order: Clostridiales\textCR%
Family: Ruminococcaceae\textCR%
Genus: Subdoligranulum\textCR%
Species: not specified%
}

\newcommand{\taxLineButyricicoccus}{%
Domain: Bacteria\textCR%
Phylum: Firmicutes\textCR%
Class: Clostridia\textCR%
Order: Clostridiales\textCR%
Family: Ruminococcaceae\textCR%
Genus: Butyricicoccus\textCR%
Species: not specified%
}

\newcommand{\taxLineLachnoclostridium}{%
Domain: Bacteria\textCR%
Phylum: Firmicutes\textCR%
Class: Clostridia\textCR%
Order: Clostridiales\textCR%
Family: Lachnospiraceae\textCR%
Genus: Lachnoclostridium\textCR%
Species: not specified%
}

\newcommand{\taxLineCoprobacter}{%
Domain: Bacteria\textCR%
Phylum: Bacteroidetes\textCR%
Class: Bacteroidia\textCR%
Order: Bacteroidales\textCR%
Family: Porphyromonadaceae\textCR%
Genus: Coprobacter\textCR%
Species: not specified%
}

\newcommand{\taxLineClostridiumInnocuumGroup}{%
Domain: Bacteria\textCR%
Phylum: Firmicutes\textCR%
Class: Erysipelotrichia\textCR%
Order: Erysipelotrichales\textCR%
Family: Erysipelotrichaceae\textCR%
Genus-level cluster: Clostridium (innocuum group)\textCR%
Species: not resolved%
}

\newcommand{\taxLineAlistipes}{%
Domain: Bacteria\textCR%
Phylum: Bacteroidetes\textCR%
Class: Bacteroidia\textCR%
Order: Bacteroidales\textCR%
Family: Rikenellaceae\textCR%
Genus: Alistipes\textCR%
Species: not specified%
}

\newcommand{\taxLineCampylobacter}{%
Domain: Bacteria\textCR%
Phylum: Proteobacteria\textCR%
Class: Epsilonproteobacteria\textCR%
Order: Campylobacterales\textCR%
Family: Campylobacteraceae\textCR%
Genus: Campylobacter\textCR%
Species: not specified%
}

\newcommand{\taxLineLachnospiraceae}{%
Domain: Bacteria\textCR%
Phylum: Firmicutes\textCR%
Class: Clostridia\textCR%
Order: Clostridiales\textCR%
Family: Lachnospiraceae%
}

\newcommand{\taxLineDorea}{%
Domain: Bacteria\textCR%
Phylum: Firmicutes\textCR%
Class: Clostridia\textCR%
Order: Clostridiales\textCR%
Family: Lachnospiraceae\textCR%
Genus: Dorea\textCR%
Species: not specified%
}

\newcommand{\taxLineOscillospira}{%
Domain: Bacteria\textCR%
Phylum: Firmicutes\textCR%
Class: Clostridia\textCR%
Order: Clostridiales\textCR%
Family: Ruminococcaceae\textCR%
Genus: Oscillospira\textCR%
Species: not specified%
}

\newcommand{\taxLineBacteroidaceae}{%
Domain: Bacteria\textCR%
Phylum: Bacteroidetes\textCR%
Class: Bacteroidia\textCR%
Order: Bacteroidales\textCR%
Family: Bacteroidaceae%
}

\newcommand{\taxLineStreptococcus}{%
Domain: Bacteria\textCR%
Phylum: Firmicutes\textCR%
Class: Bacilli\textCR%
Order: Lactobacillales\textCR%
Family: Streptococcaceae\textCR%
Genus: Streptococcus\textCR%
Species: not specified%
}

\newcommand{\taxLineEnterococcaceae}{%
Domain: Bacteria\textCR%
Phylum: Firmicutes\textCR%
Class: Bacilli\textCR%
Order: Lactobacillales\textCR%
Family: Enterococcaceae%
}

\newcommand{\taxLineLactobacillaceae}{%
Domain: Bacteria\textCR%
Phylum: Firmicutes\textCR%
Class: Bacilli\textCR%
Order: Lactobacillales\textCR%
Family: Lactobacillaceae%
}

\newcommand{\taxLineEnterobacteriaceae}{%
Domain: Bacteria\textCR%
Phylum: Proteobacteria\textCR%
Class: Gammaproteobacteria\textCR%
Order: Enterobacteriales\textCR%
Family: Enterobacteriaceae%
}

\newcommand{\taxLineStaphylococcaceae}{%
Domain: Bacteria\textCR%
Phylum: Firmicutes\textCR%
Class: Bacilli\textCR%
Order: Bacillales\textCR%
Family: Staphylococcaceae%
}

% --------------------------------------------------------------------------
% New Lineages added for Post-Transplant Multimodal Summary
% --------------------------------------------------------------------------

\newcommand{\taxLineClostridia}{%
Domain: Bacteria\textCR%
Phylum: Firmicutes\textCR%
Class: Clostridia\textCR%
Order: not specified%
}

\newcommand{\taxLineBacteroidia}{%
Domain: Bacteria\textCR%
Phylum: Bacteroidetes\textCR%
Class: Bacteroidia\textCR%
Order: not specified%
}

\newcommand{\taxLineGammaproteobacteria}{%
Domain: Bacteria\textCR%
Phylum: Proteobacteria\textCR%
Class: Gammaproteobacteria\textCR%
Order: not specified%
}

\newcommand{\taxLineStreptococcusSalivarius}{%
Domain: Bacteria\textCR%
Phylum: Firmicutes\textCR%
Class: Bacilli\textCR%
Order: Lactobacillales\textCR%
Family: Streptococcaceae\textCR%
Genus: Streptococcus\textCR%
Species: Streptococcus salivarius%
}

\newcommand{\taxLineFaecalibacteriumPrausnitzii}{%
Domain: Bacteria\textCR%
Phylum: Firmicutes\textCR%
Class: Clostridia\textCR%
Order: Clostridiales\textCR%
Family: Ruminococcaceae\textCR%
Genus: Faecalibacterium\textCR%
Species: Faecalibacterium prausnitzii%
}

\newcommand{\taxLineBlautiaWexlerae}{%
Domain: Bacteria\textCR%
Phylum: Firmicutes\textCR%
Class: Clostridia\textCR%
Order: Clostridiales\textCR%
Family: Lachnospiraceae\textCR%
Genus: Blautia\textCR%
Species: Blautia wexlerae%
}

\newcommand{\taxLineBacteroidesOvatus}{%
Domain: Bacteria\textCR%
Phylum: Bacteroidetes\textCR%
Class: Bacteroidia\textCR%
Order: Bacteroidales\textCR%
Family: Bacteroidaceae\textCR%
Genus: Bacteroides\textCR%
Species: Bacteroides ovatus%
}

% --------------------------------------------------------------------------
% Named manuscript commands
% --------------------------------------------------------------------------

% Order and family
\DeclareRobustCommand{\taxClostridiales}{%
  \taxonTooltip{\taxOrder{Clostridiales}}{Clostridiales}{\taxLineClostridiales}\xspace}

\DeclareRobustCommand{\taxRuminococcaceae}{%
  \taxonTooltip{\taxFamily{Ruminococcaceae}}{Ruminococcaceae}{\taxLineRuminococcaceae}\xspace}

\DeclareRobustCommand{\taxVeillonellaceae}{%
  \taxonTooltip{\taxFamily{Veillonellaceae}}{Veillonellaceae}{\taxLineVeillonellaceae}\xspace}

\DeclareRobustCommand{\taxSelenomonadales}{%
  \taxonTooltip{\taxOrder{Selenomonadales}}{Selenomonadales}{\taxLineSelenomonadales}\xspace}

% Genera and genus-level 16S clusters
\DeclareRobustCommand{\taxBlautia}{%
  \taxonTooltip{\taxGenus{\textit{Blautia}}}{Blautia}{\taxLineBlautia}\xspace}

\DeclareRobustCommand{\taxFaecalibacterium}{%
  \taxonTooltip{\taxGenus{\textit{Faecalibacterium}}}{Faecalibacterium}{\taxLineFaecalibacterium}\xspace}

\DeclareRobustCommand{\taxOscillibacter}{%
  \taxonTooltip{\taxGenus{\textit{Oscillibacter}}}{Oscillibacter}{\taxLineOscillibacter}\xspace}

\DeclareRobustCommand{\taxClostridiumIV}{%
  \taxonTooltip{\taxGenus{\textit{Clostridium} IV}}{Clostridium IV}{\taxLineClostridiumIV}\xspace}

\DeclareRobustCommand{\taxFlavonifractor}{%
  \taxonTooltip{\taxGenus{\textit{Flavonifractor}}}{Flavonifractor}{\taxLineFlavonifractor}\xspace}

\DeclareRobustCommand{\taxParabacteroides}{%
  \taxonTooltip{\taxGenus{\textit{Parabacteroides}}}{Parabacteroides}{\taxLineParabacteroides}\xspace}

\DeclareRobustCommand{\taxStaphylococcus}{%
  \taxonTooltip{\taxGenus{\textit{Staphylococcus}}}{Staphylococcus}{\taxLineStaphylococcus}\xspace}

\DeclareRobustCommand{\taxClostridiumXIVb}{%
  \taxonTooltip{\taxGenus{\textit{Clostridium} XIVb}}{Clostridium XIVb}{\taxLineClostridiumXIVb}\xspace}

\DeclareRobustCommand{\taxBacteroides}{%
  \taxonTooltip{\taxGenus{\textit{Bacteroides}}}{Bacteroides}{\taxLineBacteroides}\xspace}

% Species-level and unresolved species-level features
\DeclareRobustCommand{\taxAnaeromassilibacillusSp}{%
  \taxonTooltip{\taxSpecies{\textit{Anaeromassilibacillus} sp.}}{Anaeromassilibacillus sp.}{\taxLineAnaeromassilibacillusSp}\xspace}

\DeclareRobustCommand{\taxButyricicoccusSp}{%
  \taxonTooltip{\taxSpecies{\textit{Butyricicoccus} sp.}}{Butyricicoccus sp.}{\taxLineButyricicoccusSp}\xspace}

% Same shotgun feature, displayed without "sp." in the severity bullet.
\DeclareRobustCommand{\taxButyricicoccusBare}{%
  \taxonTooltip{\taxSpecies{\textit{Butyricicoccus}}}{Butyricicoccus}{\taxLineButyricicoccusSp}\xspace}

\DeclareRobustCommand{\taxLachnoanaerobaculumSpecies}{%
  \taxonTooltip{\taxSpecies{unclassified \textit{Lachnoanaerobaculum} species}}{unclassified Lachnoanaerobaculum species}{\taxLineLachnoanaerobaculumSp}\xspace}

\DeclareRobustCommand{\taxEnterococcusFaecium}{%
  \taxonTooltip{\taxSpecies{\textit{Enterococcus faecium}}}{Enterococcus faecium}{\taxLineEnterococcusFaecium}\xspace}

\DeclareRobustCommand{\taxEfaecium}{%
  \taxonTooltip{\taxSpecies{\textit{E. faecium}}}{E. faecium}{\taxLineEnterococcusFaecium}\xspace}

\DeclareRobustCommand{\taxStaphylococcusSpecies}{%
  \taxonTooltip{\taxSpecies{\textit{Staphylococcus} species}}{Staphylococcus species}{\taxLineStaphylococcusSpecies}\xspace}

\DeclareRobustCommand{\taxStaphylococcusSpp}{%
  \taxonTooltip{\taxSpecies{\textit{Staphylococcus} spp.}}{Staphylococcus spp.}{\taxLineStaphylococcusSpecies}\xspace}

\DeclareRobustCommand{\taxStaphylococcusAureus}{%
  \taxonTooltip{\taxSpecies{\textit{Staphylococcus aureus}}}{Staphylococcus aureus}{\taxLineStaphylococcusAureus}\xspace}

\DeclareRobustCommand{\taxSaureus}{%
  \taxonTooltip{\taxSpecies{\textit{S. aureus}}}{S. aureus}{\taxLineStaphylococcusAureus}\xspace}

\DeclareRobustCommand{\taxStaphylococcusEpidermidis}{%
  \taxonTooltip{\taxSpecies{\textit{Staphylococcus epidermidis}}}{Staphylococcus epidermidis}{\taxLineStaphylococcusEpidermidis}\xspace}

\DeclareRobustCommand{\taxSepidermidis}{%
  \taxonTooltip{\taxSpecies{\textit{S. epidermidis}}}{S. epidermidis}{\taxLineStaphylococcusEpidermidis}\xspace}

\DeclareRobustCommand{\taxStaphylococcusSimulans}{%
  \taxonTooltip{\taxSpecies{\textit{Staphylococcus simulans}}}{Staphylococcus simulans}{\taxLineStaphylococcusSimulans}\xspace}

\DeclareRobustCommand{\taxSsimulans}{%
  \taxonTooltip{\taxSpecies{\textit{S. simulans}}}{S. simulans}{\taxLineStaphylococcusSimulans}\xspace}

\DeclareRobustCommand{\taxBacteroidesSpecies}{%
  \taxonTooltip{\taxSpecies{\textit{Bacteroides} species}}{Bacteroides species}{\taxLineBacteroidesSpecies}\xspace}

\DeclareRobustCommand{\taxParabacteroidesDistasonis}{%
  \taxonTooltip{\taxSpecies{\textit{Parabacteroides distasonis}}}{Parabacteroides distasonis}{\taxLineParabacteroidesDistasonis}\xspace}

\DeclareRobustCommand{\taxBarnesiellaceaeSpp}{%
  \taxonTooltip{\taxFamily{Barnesiellaceae spp.}}{Barnesiellaceae spp.}{\taxLineBarnesiellaceae}\xspace}

\DeclareRobustCommand{\taxLachnobacteriumSpp}{%
  \taxonTooltip{\taxSpecies{\textit{Lachnobacterium} spp.}}{Lachnobacterium spp.}{\taxLineLachnobacterium}\xspace}

\DeclareRobustCommand{\taxHaemophilusParainfluenzae}{%
  \taxonTooltip{\taxSpecies{\textit{Haemophilus parainfluenzae}}}{Haemophilus parainfluenzae}{\taxLineHaemophilus}\xspace}

\DeclareRobustCommand{\taxEnterococcus}{%
  \taxonTooltip{\taxGenus{\textit{Enterococcus}}}{Enterococcus}{\taxLineEnterococcus}\xspace}

\DeclareRobustCommand{\taxEnterococcusFaecalis}{%
  \taxonTooltip{\taxSpecies{\textit{Enterococcus faecalis}}}{Enterococcus faecalis}{\taxLineEnterococcusFaecalis}\xspace}

\DeclareRobustCommand{\taxEfaecalis}{%
  \taxonTooltip{\taxSpecies{\textit{E. faecalis}}}{E. faecalis}{\taxLineEnterococcusFaecalis}\xspace}

\DeclareRobustCommand{\taxLactobacillus}{%
  \taxonTooltip{\taxGenus{\textit{Lactobacillus}}}{Lactobacillus}{\taxLineLactobacillus}\xspace}

\DeclareRobustCommand{\taxSubdoligranulum}{%
  \taxonTooltip{\taxGenus{\textit{Subdoligranulum}}}{Subdoligranulum}{\taxLineSubdoligranulum}\xspace}

\DeclareRobustCommand{\taxButyricicoccus}{%
  \taxonTooltip{\taxGenus{\textit{Butyricicoccus}}}{Butyricicoccus}{\taxLineButyricicoccus}\xspace}

\DeclareRobustCommand{\taxLachnoclostridium}{%
  \taxonTooltip{\taxGenus{\textit{Lachnoclostridium}}}{Lachnoclostridium}{\taxLineLachnoclostridium}\xspace}

\DeclareRobustCommand{\taxCoprobacter}{%
  \taxonTooltip{\taxGenus{\textit{Coprobacter}}}{Coprobacter}{\taxLineCoprobacter}\xspace}

\DeclareRobustCommand{\taxClostridiumInnocuumGroup}{%
  \taxonTooltip{\taxGenus{\textit{Clostridium} (innocuum group)}}{Clostridium (innocuum group)}{\taxLineClostridiumInnocuumGroup}\xspace}

\DeclareRobustCommand{\taxAlistipes}{%
  \taxonTooltip{\taxGenus{\textit{Alistipes}}}{Alistipes}{\taxLineAlistipes}\xspace}

\DeclareRobustCommand{\taxCampylobacter}{%
  \taxonTooltip{\taxGenus{\textit{Campylobacter}}}{Campylobacter}{\taxLineCampylobacter}\xspace}

\DeclareRobustCommand{\taxLachnospiraceae}{%
  \taxonTooltip{\taxFamily{Lachnospiraceae}}{Lachnospiraceae}{\taxLineLachnospiraceae}\xspace}

\DeclareRobustCommand{\taxDorea}{%
  \taxonTooltip{\taxGenus{\textit{Dorea}}}{Dorea}{\taxLineDorea}\xspace}

\DeclareRobustCommand{\taxOscillospira}{%
  \taxonTooltip{\taxGenus{\textit{Oscillospira}}}{Oscillospira}{\taxLineOscillospira}\xspace}

\DeclareRobustCommand{\taxBacteroidaceae}{%
  \taxonTooltip{\taxFamily{Bacteroidaceae}}{Bacteroidaceae}{\taxLineBacteroidaceae}\xspace}

\DeclareRobustCommand{\taxStreptococcus}{%
  \taxonTooltip{\taxGenus{\textit{Streptococcus}}}{Streptococcus}{\taxLineStreptococcus}\xspace}

\DeclareRobustCommand{\taxEnterococcaceae}{%
  \taxonTooltip{\taxFamily{Enterococcaceae}}{Enterococcaceae}{\taxLineEnterococcaceae}\xspace}

\DeclareRobustCommand{\taxLactobacillaceae}{%
  \taxonTooltip{\taxFamily{Lactobacillaceae}}{Lactobacillaceae}{\taxLineLactobacillaceae}\xspace}

\DeclareRobustCommand{\taxEnterobacteriaceae}{%
  \taxonTooltip{\taxFamily{Enterobacteriaceae}}{Enterobacteriaceae}{\taxLineEnterobacteriaceae}\xspace}

\DeclareRobustCommand{\taxStaphylococcaceae}{%
  \taxonTooltip{\taxFamily{Staphylococcaceae}}{Staphylococcaceae}{\taxLineStaphylococcaceae}\xspace}

\DeclareRobustCommand{\taxEnterococcusSpp}{%
  \taxonTooltip{\taxSpecies{\textit{Enterococcus} spp.}}{Enterococcus spp.}{\taxLineEnterococcus}\xspace}

\DeclareRobustCommand{\taxLactobacillusSpp}{%
  \taxonTooltip{\taxSpecies{\textit{Lactobacillus} spp.}}{Lactobacillus spp.}{\taxLineLactobacillus}\xspace}

% --------------------------------------------------------------------------
% New Commands added for Post-Transplant Multimodal Summary
% --------------------------------------------------------------------------

\DeclareRobustCommand{\taxClostridia}{%
  \taxonTooltip{\taxClass{Clostridia}}{Clostridia}{\taxLineClostridia}\xspace}

\DeclareRobustCommand{\taxBacteroidia}{%
  \taxonTooltip{\taxClass{Bacteroidia}}{Bacteroidia}{\taxLineBacteroidia}\xspace}

\DeclareRobustCommand{\taxGammaproteobacteria}{%
  \taxonTooltip{\taxClass{Gammaproteobacteria}}{Gammaproteobacteria}{\taxLineGammaproteobacteria}\xspace}

\DeclareRobustCommand{\taxStreptococcusSalivarius}{%
  \taxonTooltip{\taxSpecies{\textit{Streptococcus salivarius}}}{Streptococcus salivarius}{\taxLineStreptococcusSalivarius}\xspace}

\DeclareRobustCommand{\taxFaecalibacteriumPrausnitzii}{%
  \taxonTooltip{\taxSpecies{\textit{Faecalibacterium prausnitzii}}}{Faecalibacterium prausnitzii}{\taxLineFaecalibacteriumPrausnitzii}\xspace}

\DeclareRobustCommand{\taxBlautiaWexlerae}{%
  \taxonTooltip{\taxSpecies{\textit{Blautia wexlerae}}}{Blautia wexlerae}{\taxLineBlautiaWexlerae}\xspace}

\DeclareRobustCommand{\taxBacteroidesOvatus}{%
  \taxonTooltip{\taxSpecies{\textit{Bacteroides ovatus}}}{Bacteroides ovatus}{\taxLineBacteroidesOvatus}\xspace} line in the main document.

% --------------------------------------------------------------------------
% Tooltip engine
% --------------------------------------------------------------------------
\newif\ifTaxTooltips
\TaxTooltipstrue

% Arguments:
%   #1 = formatted text printed in the manuscript
%   #2 = plain-text fallback for PDF bookmarks
%   #3 = tooltip text
\DeclareRobustCommand{\taxonTooltip}[3]{%
  \texorpdfstring{%
    \ifTaxTooltips
      \pdftooltip{#1}{#3}%
    \else
      #1%
    \fi
  }{#2}%
}

% --------------------------------------------------------------------------
% Reusable lineage strings
% --------------------------------------------------------------------------
\newcommand{\taxLineClostridiales}{%
Domain: Bacteria\textCR%
Phylum: Firmicutes\textCR%
Class: Clostridia\textCR%
Order: Clostridiales%
}

\newcommand{\taxLineRuminococcaceae}{%
Domain: Bacteria\textCR%
Phylum: Firmicutes\textCR%
Class: Clostridia\textCR%
Order: Clostridiales\textCR%
Family: Ruminococcaceae%
}

\newcommand{\taxLineBlautia}{%
Domain: Bacteria\textCR%
Phylum: Firmicutes\textCR%
Class: Clostridia\textCR%
Order: Clostridiales\textCR%
Family: Lachnospiraceae\textCR%
Genus: Blautia\textCR%
Species: not specified%
}

\newcommand{\taxLineFaecalibacterium}{%
Domain: Bacteria\textCR%
Phylum: Firmicutes\textCR%
Class: Clostridia\textCR%
Order: Clostridiales\textCR%
Family: Ruminococcaceae\textCR%
Genus: Faecalibacterium\textCR%
Species: not specified%
}

\newcommand{\taxLineOscillibacter}{%
Domain: Bacteria\textCR%
Phylum: Firmicutes\textCR%
Class: Clostridia\textCR%
Order: Clostridiales\textCR%
Family: Ruminococcaceae\textCR%
Genus: Oscillibacter\textCR%
Species: not specified%
}

\newcommand{\taxLineClostridiumIV}{%
Domain: Bacteria\textCR%
Phylum: Firmicutes\textCR%
Class: Clostridia\textCR%
Order: Clostridiales\textCR%
Family: Ruminococcaceae\textCR%
Genus-level 16S cluster: Clostridium IV\textCR%
Species: not resolved%
}

\newcommand{\taxLineFlavonifractor}{%
Domain: Bacteria\textCR%
Phylum: Firmicutes\textCR%
Class: Clostridia\textCR%
Order: Clostridiales\textCR%
Family: Ruminococcaceae\textCR%
Genus: Flavonifractor\textCR%
Species: not specified%
}

\newcommand{\taxLineAnaeromassilibacillusSp}{%
Domain: Bacteria\textCR%
Phylum: Firmicutes\textCR%
Class: Clostridia\textCR%
Order: Clostridiales\textCR%
Family: Ruminococcaceae\textCR%
Genus: Anaeromassilibacillus\textCR%
Species: unresolved (reported as Anaeromassilibacillus sp.)%
}

\newcommand{\taxLineButyricicoccusSp}{%
Domain: Bacteria\textCR%
Phylum: Firmicutes\textCR%
Class: Clostridia\textCR%
Order: Clostridiales\textCR%
Family: Ruminococcaceae\textCR%
Genus: Butyricicoccus\textCR%
Species: unresolved (reported as Butyricicoccus sp.)%
}

\newcommand{\taxLineLachnoanaerobaculumSp}{%
Domain: Bacteria\textCR%
Phylum: Firmicutes\textCR%
Class: Clostridia\textCR%
Order: Clostridiales\textCR%
Family: Lachnospiraceae\textCR%
Genus: Lachnoanaerobaculum\textCR%
Species: unresolved/unclassified%
}

\newcommand{\taxLineEnterococcusFaecium}{%
Domain: Bacteria\textCR%
Phylum: Firmicutes\textCR%
Class: Bacilli\textCR%
Order: Lactobacillales\textCR%
Family: Enterococcaceae\textCR%
Genus: Enterococcus\textCR%
Species: Enterococcus faecium%
}

\newcommand{\taxLineParabacteroides}{%
Domain: Bacteria\textCR%
Phylum: Bacteroidetes\textCR%
Class: Bacteroidia\textCR%
Order: Bacteroidales\textCR%
Family: Tannerellaceae\textCR%
Genus: Parabacteroides\textCR%
Species: not specified%
}

\newcommand{\taxLineStaphylococcus}{%
Domain: Bacteria\textCR%
Phylum: Firmicutes\textCR%
Class: Bacilli\textCR%
Order: Bacillales\textCR%
Family: Staphylococcaceae\textCR%
Genus: Staphylococcus\textCR%
Species: not specified%
}

\newcommand{\taxLineStaphylococcusSpecies}{%
Domain: Bacteria\textCR%
Phylum: Firmicutes\textCR%
Class: Bacilli\textCR%
Order: Bacillales\textCR%
Family: Staphylococcaceae\textCR%
Genus: Staphylococcus\textCR%
Species discussed: S. aureus; S. epidermidis; S. simulans%
}

\newcommand{\taxLineStaphylococcusAureus}{%
Domain: Bacteria\textCR%
Phylum: Firmicutes\textCR%
Class: Bacilli\textCR%
Order: Bacillales\textCR%
Family: Staphylococcaceae\textCR%
Genus: Staphylococcus\textCR%
Species: Staphylococcus aureus%
}

\newcommand{\taxLineStaphylococcusEpidermidis}{%
Domain: Bacteria\textCR%
Phylum: Firmicutes\textCR%
Class: Bacilli\textCR%
Order: Bacillales\textCR%
Family: Staphylococcaceae\textCR%
Genus: Staphylococcus\textCR%
Species: Staphylococcus epidermidis%
}

\newcommand{\taxLineStaphylococcusSimulans}{%
Domain: Bacteria\textCR%
Phylum: Firmicutes\textCR%
Class: Bacilli\textCR%
Order: Bacillales\textCR%
Family: Staphylococcaceae\textCR%
Genus: Staphylococcus\textCR%
Species: Staphylococcus simulans%
}

\newcommand{\taxLineClostridiumXIVb}{%
Domain: Bacteria\textCR%
Phylum: Firmicutes\textCR%
Class: Clostridia\textCR%
Order: Clostridiales\textCR%
Family: not resolved by this paper-level label\textCR%
Genus-level 16S cluster: Clostridium XIVb\textCR%
Species: not resolved%
}

\newcommand{\taxLineBacteroides}{%
Domain: Bacteria\textCR%
Phylum: Bacteroidetes\textCR%
Class: Bacteroidia\textCR%
Order: Bacteroidales\textCR%
Family: Bacteroidaceae\textCR%
Genus: Bacteroides\textCR%
Species: not specified%
}

\newcommand{\taxLineBacteroidesSpecies}{%
Domain: Bacteria\textCR%
Phylum: Bacteroidetes\textCR%
Class: Bacteroidia\textCR%
Order: Bacteroidales\textCR%
Family: Bacteroidaceae\textCR%
Genus: Bacteroides\textCR%
Species: unspecified/multiple species%
}

\newcommand{\taxLineVeillonellaceae}{%
Domain: Bacteria\textCR%
Phylum: Firmicutes\textCR%
Class: Negativicutes\textCR%
Order: Selenomonadales\textCR%
Family: Veillonellaceae%
}

\newcommand{\taxLineSelenomonadales}{%
Domain: Bacteria\textCR%
Phylum: Firmicutes\textCR%
Class: Negativicutes\textCR%
Order: Selenomonadales%
}

\newcommand{\taxLineParabacteroidesDistasonis}{%
Domain: Bacteria\textCR%
Phylum: Bacteroidetes\textCR%
Class: Bacteroidia\textCR%
Order: Bacteroidales\textCR%
Family: Porphyromonadaceae\textCR%
Genus: Parabacteroides\textCR%
Species: Parabacteroides distasonis%
}

\newcommand{\taxLineBarnesiellaceae}{%
Domain: Bacteria\textCR%
Phylum: Bacteroidetes\textCR%
Class: Bacteroidia\textCR%
Order: Bacteroidales\textCR%
Family: Barnesiellaceae%
}

\newcommand{\taxLineLachnobacterium}{%
Domain: Bacteria\textCR%
Phylum: Firmicutes\textCR%
Class: Clostridia\textCR%
Order: Clostridiales\textCR%
Family: Lachnospiraceae\textCR%
Genus: Lachnobacterium\textCR%
Species: unresolved (reported as Lachnobacterium unknown)%
}

\newcommand{\taxLineHaemophilus}{%
Domain: Bacteria\textCR%
Phylum: Proteobacteria\textCR%
Class: Gammaproteobacteria\textCR%
Order: Pasteurellales\textCR%
Family: Pasteurellaceae\textCR%
Genus: Haemophilus\textCR%
Species: unresolved (reported as Haemophilus unknown)%
}

\newcommand{\taxLineEnterococcus}{%
Domain: Bacteria\textCR%
Phylum: Firmicutes\textCR%
Class: Bacilli\textCR%
Order: Lactobacillales\textCR%
Family: Enterococcaceae\textCR%
Genus: Enterococcus\textCR%
Species: not specified%
}

\newcommand{\taxLineEnterococcusFaecalis}{%
Domain: Bacteria\textCR%
Phylum: Firmicutes\textCR%
Class: Bacilli\textCR%
Order: Lactobacillales\textCR%
Family: Enterococcaceae\textCR%
Genus: Enterococcus\textCR%
Species: Enterococcus faecalis%
}

\newcommand{\taxLineLactobacillus}{%
Domain: Bacteria\textCR%
Phylum: Firmicutes\textCR%
Class: Bacilli\textCR%
Order: Lactobacillales\textCR%
Family: Lactobacillaceae\textCR%
Genus: Lactobacillus\textCR%
Species: not specified%
}

\newcommand{\taxLineSubdoligranulum}{%
Domain: Bacteria\textCR%
Phylum: Firmicutes\textCR%
Class: Clostridia\textCR%
Order: Clostridiales\textCR%
Family: Ruminococcaceae\textCR%
Genus: Subdoligranulum\textCR%
Species: not specified%
}

\newcommand{\taxLineButyricicoccus}{%
Domain: Bacteria\textCR%
Phylum: Firmicutes\textCR%
Class: Clostridia\textCR%
Order: Clostridiales\textCR%
Family: Ruminococcaceae\textCR%
Genus: Butyricicoccus\textCR%
Species: not specified%
}

\newcommand{\taxLineLachnoclostridium}{%
Domain: Bacteria\textCR%
Phylum: Firmicutes\textCR%
Class: Clostridia\textCR%
Order: Clostridiales\textCR%
Family: Lachnospiraceae\textCR%
Genus: Lachnoclostridium\textCR%
Species: not specified%
}

\newcommand{\taxLineCoprobacter}{%
Domain: Bacteria\textCR%
Phylum: Bacteroidetes\textCR%
Class: Bacteroidia\textCR%
Order: Bacteroidales\textCR%
Family: Porphyromonadaceae\textCR%
Genus: Coprobacter\textCR%
Species: not specified%
}

\newcommand{\taxLineClostridiumInnocuumGroup}{%
Domain: Bacteria\textCR%
Phylum: Firmicutes\textCR%
Class: Erysipelotrichia\textCR%
Order: Erysipelotrichales\textCR%
Family: Erysipelotrichaceae\textCR%
Genus-level cluster: Clostridium (innocuum group)\textCR%
Species: not resolved%
}

\newcommand{\taxLineAlistipes}{%
Domain: Bacteria\textCR%
Phylum: Bacteroidetes\textCR%
Class: Bacteroidia\textCR%
Order: Bacteroidales\textCR%
Family: Rikenellaceae\textCR%
Genus: Alistipes\textCR%
Species: not specified%
}

\newcommand{\taxLineCampylobacter}{%
Domain: Bacteria\textCR%
Phylum: Proteobacteria\textCR%
Class: Epsilonproteobacteria\textCR%
Order: Campylobacterales\textCR%
Family: Campylobacteraceae\textCR%
Genus: Campylobacter\textCR%
Species: not specified%
}

\newcommand{\taxLineLachnospiraceae}{%
Domain: Bacteria\textCR%
Phylum: Firmicutes\textCR%
Class: Clostridia\textCR%
Order: Clostridiales\textCR%
Family: Lachnospiraceae%
}

\newcommand{\taxLineDorea}{%
Domain: Bacteria\textCR%
Phylum: Firmicutes\textCR%
Class: Clostridia\textCR%
Order: Clostridiales\textCR%
Family: Lachnospiraceae\textCR%
Genus: Dorea\textCR%
Species: not specified%
}

\newcommand{\taxLineOscillospira}{%
Domain: Bacteria\textCR%
Phylum: Firmicutes\textCR%
Class: Clostridia\textCR%
Order: Clostridiales\textCR%
Family: Ruminococcaceae\textCR%
Genus: Oscillospira\textCR%
Species: not specified%
}

\newcommand{\taxLineBacteroidaceae}{%
Domain: Bacteria\textCR%
Phylum: Bacteroidetes\textCR%
Class: Bacteroidia\textCR%
Order: Bacteroidales\textCR%
Family: Bacteroidaceae%
}

\newcommand{\taxLineStreptococcus}{%
Domain: Bacteria\textCR%
Phylum: Firmicutes\textCR%
Class: Bacilli\textCR%
Order: Lactobacillales\textCR%
Family: Streptococcaceae\textCR%
Genus: Streptococcus\textCR%
Species: not specified%
}

\newcommand{\taxLineEnterococcaceae}{%
Domain: Bacteria\textCR%
Phylum: Firmicutes\textCR%
Class: Bacilli\textCR%
Order: Lactobacillales\textCR%
Family: Enterococcaceae%
}

\newcommand{\taxLineLactobacillaceae}{%
Domain: Bacteria\textCR%
Phylum: Firmicutes\textCR%
Class: Bacilli\textCR%
Order: Lactobacillales\textCR%
Family: Lactobacillaceae%
}

\newcommand{\taxLineEnterobacteriaceae}{%
Domain: Bacteria\textCR%
Phylum: Proteobacteria\textCR%
Class: Gammaproteobacteria\textCR%
Order: Enterobacteriales\textCR%
Family: Enterobacteriaceae%
}

\newcommand{\taxLineStaphylococcaceae}{%
Domain: Bacteria\textCR%
Phylum: Firmicutes\textCR%
Class: Bacilli\textCR%
Order: Bacillales\textCR%
Family: Staphylococcaceae%
}

% --------------------------------------------------------------------------
% New Lineages added for Post-Transplant Multimodal Summary
% --------------------------------------------------------------------------

\newcommand{\taxLineClostridia}{%
Domain: Bacteria\textCR%
Phylum: Firmicutes\textCR%
Class: Clostridia\textCR%
Order: not specified%
}

\newcommand{\taxLineBacteroidia}{%
Domain: Bacteria\textCR%
Phylum: Bacteroidetes\textCR%
Class: Bacteroidia\textCR%
Order: not specified%
}

\newcommand{\taxLineGammaproteobacteria}{%
Domain: Bacteria\textCR%
Phylum: Proteobacteria\textCR%
Class: Gammaproteobacteria\textCR%
Order: not specified%
}

\newcommand{\taxLineStreptococcusSalivarius}{%
Domain: Bacteria\textCR%
Phylum: Firmicutes\textCR%
Class: Bacilli\textCR%
Order: Lactobacillales\textCR%
Family: Streptococcaceae\textCR%
Genus: Streptococcus\textCR%
Species: Streptococcus salivarius%
}

\newcommand{\taxLineFaecalibacteriumPrausnitzii}{%
Domain: Bacteria\textCR%
Phylum: Firmicutes\textCR%
Class: Clostridia\textCR%
Order: Clostridiales\textCR%
Family: Ruminococcaceae\textCR%
Genus: Faecalibacterium\textCR%
Species: Faecalibacterium prausnitzii%
}

\newcommand{\taxLineBlautiaWexlerae}{%
Domain: Bacteria\textCR%
Phylum: Firmicutes\textCR%
Class: Clostridia\textCR%
Order: Clostridiales\textCR%
Family: Lachnospiraceae\textCR%
Genus: Blautia\textCR%
Species: Blautia wexlerae%
}

\newcommand{\taxLineBacteroidesOvatus}{%
Domain: Bacteria\textCR%
Phylum: Bacteroidetes\textCR%
Class: Bacteroidia\textCR%
Order: Bacteroidales\textCR%
Family: Bacteroidaceae\textCR%
Genus: Bacteroides\textCR%
Species: Bacteroides ovatus%
}

% --------------------------------------------------------------------------
% Named manuscript commands
% --------------------------------------------------------------------------

% Order and family
\DeclareRobustCommand{\taxClostridiales}{%
  \taxonTooltip{\taxOrder{Clostridiales}}{Clostridiales}{\taxLineClostridiales}\xspace}

\DeclareRobustCommand{\taxRuminococcaceae}{%
  \taxonTooltip{\taxFamily{Ruminococcaceae}}{Ruminococcaceae}{\taxLineRuminococcaceae}\xspace}

\DeclareRobustCommand{\taxVeillonellaceae}{%
  \taxonTooltip{\taxFamily{Veillonellaceae}}{Veillonellaceae}{\taxLineVeillonellaceae}\xspace}

\DeclareRobustCommand{\taxSelenomonadales}{%
  \taxonTooltip{\taxOrder{Selenomonadales}}{Selenomonadales}{\taxLineSelenomonadales}\xspace}

% Genera and genus-level 16S clusters
\DeclareRobustCommand{\taxBlautia}{%
  \taxonTooltip{\taxGenus{\textit{Blautia}}}{Blautia}{\taxLineBlautia}\xspace}

\DeclareRobustCommand{\taxFaecalibacterium}{%
  \taxonTooltip{\taxGenus{\textit{Faecalibacterium}}}{Faecalibacterium}{\taxLineFaecalibacterium}\xspace}

\DeclareRobustCommand{\taxOscillibacter}{%
  \taxonTooltip{\taxGenus{\textit{Oscillibacter}}}{Oscillibacter}{\taxLineOscillibacter}\xspace}

\DeclareRobustCommand{\taxClostridiumIV}{%
  \taxonTooltip{\taxGenus{\textit{Clostridium} IV}}{Clostridium IV}{\taxLineClostridiumIV}\xspace}

\DeclareRobustCommand{\taxFlavonifractor}{%
  \taxonTooltip{\taxGenus{\textit{Flavonifractor}}}{Flavonifractor}{\taxLineFlavonifractor}\xspace}

\DeclareRobustCommand{\taxParabacteroides}{%
  \taxonTooltip{\taxGenus{\textit{Parabacteroides}}}{Parabacteroides}{\taxLineParabacteroides}\xspace}

\DeclareRobustCommand{\taxStaphylococcus}{%
  \taxonTooltip{\taxGenus{\textit{Staphylococcus}}}{Staphylococcus}{\taxLineStaphylococcus}\xspace}

\DeclareRobustCommand{\taxClostridiumXIVb}{%
  \taxonTooltip{\taxGenus{\textit{Clostridium} XIVb}}{Clostridium XIVb}{\taxLineClostridiumXIVb}\xspace}

\DeclareRobustCommand{\taxBacteroides}{%
  \taxonTooltip{\taxGenus{\textit{Bacteroides}}}{Bacteroides}{\taxLineBacteroides}\xspace}

% Species-level and unresolved species-level features
\DeclareRobustCommand{\taxAnaeromassilibacillusSp}{%
  \taxonTooltip{\taxSpecies{\textit{Anaeromassilibacillus} sp.}}{Anaeromassilibacillus sp.}{\taxLineAnaeromassilibacillusSp}\xspace}

\DeclareRobustCommand{\taxButyricicoccusSp}{%
  \taxonTooltip{\taxSpecies{\textit{Butyricicoccus} sp.}}{Butyricicoccus sp.}{\taxLineButyricicoccusSp}\xspace}

% Same shotgun feature, displayed without "sp." in the severity bullet.
\DeclareRobustCommand{\taxButyricicoccusBare}{%
  \taxonTooltip{\taxSpecies{\textit{Butyricicoccus}}}{Butyricicoccus}{\taxLineButyricicoccusSp}\xspace}

\DeclareRobustCommand{\taxLachnoanaerobaculumSpecies}{%
  \taxonTooltip{\taxSpecies{unclassified \textit{Lachnoanaerobaculum} species}}{unclassified Lachnoanaerobaculum species}{\taxLineLachnoanaerobaculumSp}\xspace}

\DeclareRobustCommand{\taxEnterococcusFaecium}{%
  \taxonTooltip{\taxSpecies{\textit{Enterococcus faecium}}}{Enterococcus faecium}{\taxLineEnterococcusFaecium}\xspace}

\DeclareRobustCommand{\taxEfaecium}{%
  \taxonTooltip{\taxSpecies{\textit{E. faecium}}}{E. faecium}{\taxLineEnterococcusFaecium}\xspace}

\DeclareRobustCommand{\taxStaphylococcusSpecies}{%
  \taxonTooltip{\taxSpecies{\textit{Staphylococcus} species}}{Staphylococcus species}{\taxLineStaphylococcusSpecies}\xspace}

\DeclareRobustCommand{\taxStaphylococcusSpp}{%
  \taxonTooltip{\taxSpecies{\textit{Staphylococcus} spp.}}{Staphylococcus spp.}{\taxLineStaphylococcusSpecies}\xspace}

\DeclareRobustCommand{\taxStaphylococcusAureus}{%
  \taxonTooltip{\taxSpecies{\textit{Staphylococcus aureus}}}{Staphylococcus aureus}{\taxLineStaphylococcusAureus}\xspace}

\DeclareRobustCommand{\taxSaureus}{%
  \taxonTooltip{\taxSpecies{\textit{S. aureus}}}{S. aureus}{\taxLineStaphylococcusAureus}\xspace}

\DeclareRobustCommand{\taxStaphylococcusEpidermidis}{%
  \taxonTooltip{\taxSpecies{\textit{Staphylococcus epidermidis}}}{Staphylococcus epidermidis}{\taxLineStaphylococcusEpidermidis}\xspace}

\DeclareRobustCommand{\taxSepidermidis}{%
  \taxonTooltip{\taxSpecies{\textit{S. epidermidis}}}{S. epidermidis}{\taxLineStaphylococcusEpidermidis}\xspace}

\DeclareRobustCommand{\taxStaphylococcusSimulans}{%
  \taxonTooltip{\taxSpecies{\textit{Staphylococcus simulans}}}{Staphylococcus simulans}{\taxLineStaphylococcusSimulans}\xspace}

\DeclareRobustCommand{\taxSsimulans}{%
  \taxonTooltip{\taxSpecies{\textit{S. simulans}}}{S. simulans}{\taxLineStaphylococcusSimulans}\xspace}

\DeclareRobustCommand{\taxBacteroidesSpecies}{%
  \taxonTooltip{\taxSpecies{\textit{Bacteroides} species}}{Bacteroides species}{\taxLineBacteroidesSpecies}\xspace}

\DeclareRobustCommand{\taxParabacteroidesDistasonis}{%
  \taxonTooltip{\taxSpecies{\textit{Parabacteroides distasonis}}}{Parabacteroides distasonis}{\taxLineParabacteroidesDistasonis}\xspace}

\DeclareRobustCommand{\taxBarnesiellaceaeSpp}{%
  \taxonTooltip{\taxFamily{Barnesiellaceae spp.}}{Barnesiellaceae spp.}{\taxLineBarnesiellaceae}\xspace}

\DeclareRobustCommand{\taxLachnobacteriumSpp}{%
  \taxonTooltip{\taxSpecies{\textit{Lachnobacterium} spp.}}{Lachnobacterium spp.}{\taxLineLachnobacterium}\xspace}

\DeclareRobustCommand{\taxHaemophilusParainfluenzae}{%
  \taxonTooltip{\taxSpecies{\textit{Haemophilus parainfluenzae}}}{Haemophilus parainfluenzae}{\taxLineHaemophilus}\xspace}

\DeclareRobustCommand{\taxEnterococcus}{%
  \taxonTooltip{\taxGenus{\textit{Enterococcus}}}{Enterococcus}{\taxLineEnterococcus}\xspace}

\DeclareRobustCommand{\taxEnterococcusFaecalis}{%
  \taxonTooltip{\taxSpecies{\textit{Enterococcus faecalis}}}{Enterococcus faecalis}{\taxLineEnterococcusFaecalis}\xspace}

\DeclareRobustCommand{\taxEfaecalis}{%
  \taxonTooltip{\taxSpecies{\textit{E. faecalis}}}{E. faecalis}{\taxLineEnterococcusFaecalis}\xspace}

\DeclareRobustCommand{\taxLactobacillus}{%
  \taxonTooltip{\taxGenus{\textit{Lactobacillus}}}{Lactobacillus}{\taxLineLactobacillus}\xspace}

\DeclareRobustCommand{\taxSubdoligranulum}{%
  \taxonTooltip{\taxGenus{\textit{Subdoligranulum}}}{Subdoligranulum}{\taxLineSubdoligranulum}\xspace}

\DeclareRobustCommand{\taxButyricicoccus}{%
  \taxonTooltip{\taxGenus{\textit{Butyricicoccus}}}{Butyricicoccus}{\taxLineButyricicoccus}\xspace}

\DeclareRobustCommand{\taxLachnoclostridium}{%
  \taxonTooltip{\taxGenus{\textit{Lachnoclostridium}}}{Lachnoclostridium}{\taxLineLachnoclostridium}\xspace}

\DeclareRobustCommand{\taxCoprobacter}{%
  \taxonTooltip{\taxGenus{\textit{Coprobacter}}}{Coprobacter}{\taxLineCoprobacter}\xspace}

\DeclareRobustCommand{\taxClostridiumInnocuumGroup}{%
  \taxonTooltip{\taxGenus{\textit{Clostridium} (innocuum group)}}{Clostridium (innocuum group)}{\taxLineClostridiumInnocuumGroup}\xspace}

\DeclareRobustCommand{\taxAlistipes}{%
  \taxonTooltip{\taxGenus{\textit{Alistipes}}}{Alistipes}{\taxLineAlistipes}\xspace}

\DeclareRobustCommand{\taxCampylobacter}{%
  \taxonTooltip{\taxGenus{\textit{Campylobacter}}}{Campylobacter}{\taxLineCampylobacter}\xspace}

\DeclareRobustCommand{\taxLachnospiraceae}{%
  \taxonTooltip{\taxFamily{Lachnospiraceae}}{Lachnospiraceae}{\taxLineLachnospiraceae}\xspace}

\DeclareRobustCommand{\taxDorea}{%
  \taxonTooltip{\taxGenus{\textit{Dorea}}}{Dorea}{\taxLineDorea}\xspace}

\DeclareRobustCommand{\taxOscillospira}{%
  \taxonTooltip{\taxGenus{\textit{Oscillospira}}}{Oscillospira}{\taxLineOscillospira}\xspace}

\DeclareRobustCommand{\taxBacteroidaceae}{%
  \taxonTooltip{\taxFamily{Bacteroidaceae}}{Bacteroidaceae}{\taxLineBacteroidaceae}\xspace}

\DeclareRobustCommand{\taxStreptococcus}{%
  \taxonTooltip{\taxGenus{\textit{Streptococcus}}}{Streptococcus}{\taxLineStreptococcus}\xspace}

\DeclareRobustCommand{\taxEnterococcaceae}{%
  \taxonTooltip{\taxFamily{Enterococcaceae}}{Enterococcaceae}{\taxLineEnterococcaceae}\xspace}

\DeclareRobustCommand{\taxLactobacillaceae}{%
  \taxonTooltip{\taxFamily{Lactobacillaceae}}{Lactobacillaceae}{\taxLineLactobacillaceae}\xspace}

\DeclareRobustCommand{\taxEnterobacteriaceae}{%
  \taxonTooltip{\taxFamily{Enterobacteriaceae}}{Enterobacteriaceae}{\taxLineEnterobacteriaceae}\xspace}

\DeclareRobustCommand{\taxStaphylococcaceae}{%
  \taxonTooltip{\taxFamily{Staphylococcaceae}}{Staphylococcaceae}{\taxLineStaphylococcaceae}\xspace}

\DeclareRobustCommand{\taxEnterococcusSpp}{%
  \taxonTooltip{\taxSpecies{\textit{Enterococcus} spp.}}{Enterococcus spp.}{\taxLineEnterococcus}\xspace}

\DeclareRobustCommand{\taxLactobacillusSpp}{%
  \taxonTooltip{\taxSpecies{\textit{Lactobacillus} spp.}}{Lactobacillus spp.}{\taxLineLactobacillus}\xspace}

% --------------------------------------------------------------------------
% New Commands added for Post-Transplant Multimodal Summary
% --------------------------------------------------------------------------

\DeclareRobustCommand{\taxClostridia}{%
  \taxonTooltip{\taxClass{Clostridia}}{Clostridia}{\taxLineClostridia}\xspace}

\DeclareRobustCommand{\taxBacteroidia}{%
  \taxonTooltip{\taxClass{Bacteroidia}}{Bacteroidia}{\taxLineBacteroidia}\xspace}

\DeclareRobustCommand{\taxGammaproteobacteria}{%
  \taxonTooltip{\taxClass{Gammaproteobacteria}}{Gammaproteobacteria}{\taxLineGammaproteobacteria}\xspace}

\DeclareRobustCommand{\taxStreptococcusSalivarius}{%
  \taxonTooltip{\taxSpecies{\textit{Streptococcus salivarius}}}{Streptococcus salivarius}{\taxLineStreptococcusSalivarius}\xspace}

\DeclareRobustCommand{\taxFaecalibacteriumPrausnitzii}{%
  \taxonTooltip{\taxSpecies{\textit{Faecalibacterium prausnitzii}}}{Faecalibacterium prausnitzii}{\taxLineFaecalibacteriumPrausnitzii}\xspace}

\DeclareRobustCommand{\taxBlautiaWexlerae}{%
  \taxonTooltip{\taxSpecies{\textit{Blautia wexlerae}}}{Blautia wexlerae}{\taxLineBlautiaWexlerae}\xspace}

\DeclareRobustCommand{\taxBacteroidesOvatus}{%
  \taxonTooltip{\taxSpecies{\textit{Bacteroides ovatus}}}{Bacteroides ovatus}{\taxLineBacteroidesOvatus}\xspace}

% Uncomment the following line to disable every taxonomic tooltip
% while retaining all colors and italics:
% \TaxTooltipsfalse



%%%%%%%%%%%%%%%%%%%%%%%%%%%%%%%%%%%%%%%%%%%%%%%%%%%%%%%%%%%%%%%%%%%
\begin{document}
%
\title{Gut Microbiome in Adult Allo-SCT Patients: Analysis Sub Cohort}
\author{Sophie Huebler \thanks{\href{mailto:sophie.huebler@utah.edu}{sophie.huebler@utah.edu}}}
\date{\today}
\maketitle
%\vfill
\renewcommand\abstractname{Objective}
\section*{\abstractname}
\textit{Make a list of adult allo-HSCT recipient cohorts with publicly available microbiome data (and associated clinical data).}
\thispagestyle{empty}
%
%
\tableofcontents
\thispagestyle{empty}
\cleardoublepage
\pagenumbering{arabic} 
\newpage
%
%
\section{Studies}

% =====================================================================
% TABLE 1 — Study-level summary
% Columns: Reference | Microbiome Cohort | Study design | Main results
% Reference cell (rightmost) stacks: cohort name + \cite, BioProject underneath.
% Data cells left as \tbd — fill in paper-by-paper.
% =====================================================================
\begin{longtblr}[
  caption = {Human studies of the association between the gut microbiota and GVHD with publicly available 16S sequencing and clinical metadata included in the analytic sub-cohort.},
  label = {tab:study_summary},
]{
  width = \linewidth,
  colspec = {Q[l,m,0.17\linewidth] Q[l,m,0.15\linewidth] Q[l,m,0.30\linewidth] Q[l,m,0.32\linewidth] },
  row{1} = {font=\bfseries}, % Bold header row
  hline{1,2,Z} = {0.8pt},    % Rules at top, below header, and bottom
  cells = {font=\small},
  column{4} = {font=\footnotesize}, % Results column smaller
  rowhead = 1
}

Reference & Microbiome Cohort & Study design & Main results  \\
\refcell{Artacho2024}{artacho2024multimodal-9e5}{PRJNA1088414} & 105 adults undergoing allo-SCT & Longitudinal multi-omics cohort comparing stool collected shortly before stem-cell infusion and around engraftment, before GVHD onset, to identify taxonomic, functional, metabolomic, and bacterial-metabolite changes associated with allo-HSCT and subsequent moderate-to-severe acute GVHD. & Allo-HSCT was followed by reduced microbiota diversity, depletion of Clostridiales and short-chain fatty acids, and expansion of \textit{Staphylococcus}. Patients who later developed moderate-to-severe acute GVHD showed greater depletion of specific Clostridiales taxa, greater expansion of \textit{Enterococcus faecium} and \textit{Parabacteroides}, and more pronounced disruption of microbiome-encoded functions, metabolites, and bacterial-metabolite interactions.  \\
\refcell{DAmico2019}{.}{PRJNA592853} & 20 children undergoing allo-SCT & Longitudinal cohort examining whether enteral or parenteral nutrition affects gut microbiome recovery after allo-SCT.  & Both nutritional groups experienced early microbiome disruption, but only patients receiving enteral nutrition recovered pre-HSCT-like diversity, community composition, and short-chain fatty acid production. Parenteral nutrition was associated with persistent dysbiosis and all documented bloodstream infections. Acute GVHD occurred at a similar frequency in both groups, and the study did not establish that enteral nutrition reduced aGVHD incidence. \\ % TODO: add DAmico2019 to references.bib
\refcell{Fujimoto2024}{Fujimoto2024}{PRJNA1109929} & 45 adults undergoing allo-SCT & Prospective longitudinal observational cohort with serial stool 16S profiling, enterococcal culture, and isolate whole-genome and prophage analysis, followed by recombinant production and in vitro and gnotobiotic mouse testing of an \textit{E. faecalis}-specific phage-derived endolysin.  & \textit{Enterococcus} domination was common during allo-HCT, and fecal cytolysin positivity among Enterococcus-dominated cases was associated with aGVHD. A phage-derived endolysin selectively targeted biofilm-forming \textit{E. faecalis} and improved survival in preclinical aGVHD mouse models. \\
\refcell{Ingham2019}{ingham2019specific-2d2}{PRJEB25221} & 30 children undergoing allo-SCT & Prospective longitudinal cohort investigating exploratory high dimensional clustering of the gut microbiome, immune markers, and clinical outcomes.    &  Severe GVHD and death were associated with a lack of commensal bacteria like \textit{Clostridiales}, high plasma hBD2 levels post-transplant, and increased \textit{Lactobacillaceae} at all time points. NK and B cell reconstitution was associated with mild to no GVHD and better survival, and was associated with high abundance of obligate anaerobes.  \\
\refcell{Jarosch2023}{jarosch2023multimodal-0f7}{PRJEB60178} & 46 adults who underwent allo-SCT & Longitudinal multi-modal analysis relating the gut microbiome, immune cell function, and GVHD.  & Broad-spectrum antibiotic exposure was associated with lower microbial diversity and fewer, less suppressive intestinal Tregs. SCFA-producing taxa, particularly Clostridia and Bacteroidia, were associated with suppressive Treg phenotypes, while severe GI-GVHD was associated with persistent, clonally expanded activated CD8 T cells. Overall microbial community composition did not separate according to GI-GVHD severity.  \\
\refcell{Liu2017}{Liu2017}{PRJEB16057} & 57 adults undergoing allo-SCT; 22 healthy SCT donors & Prospective cohort examining allo-SCT recipient and paired sibling-donor gut microbiota, at the time of pre-conditioning, in relation to GI-GVHD and mortality. & Recipients had lower baseline diversity and different composition than healthy donors. Lower recipient diversity was associated with higher overall mortality but not GI-GVHD. Donor-recipient dissimilarity was not associated with GI-GVHD, whereas higher donor diversity was associated with lower GI-GVHD odds.  \\
\refcell{Spindelboeck2021}{.}{ PRJEB39834} & 10 adults with GI-GVHD after allo-SCT, 14 healthy FMT donors & Prospective cohort of adults undergoing repeated colonoscopic FMT for severe treatment-refractory GI-GVHD after allo-SCT. &  Patients who responded to FMT with full or partial GI-GVHD remission at 90 days had longer overall and GI-GVHD-related survival. Responders developed a more donor-like microbiota, with restored diversity and enrichment of anaerobic butyrate-producing taxa. Ongoing broad-spectrum antibiotics were associated with persistent dysbiosis, failed engraftment, and non-response to FMT.  \\
\refcell{Vallet2023}{.}{PRJNA902819}  & 55 adults undergoing allo-SCT  & Longidudinal multi-omics substudy of an RCT where stool samples were used to form bacterial enterotypes which interacted with viral and metabolic pathways and were investigated for association with receipt of the drug and allo-SCT outcomes.   & One specific enterotype was associated with greater odds of remission, and changing enterotypes generally was associated with greater odds of remission. Enterotype distribution was not associated with acute or chronic GVHD. \\ % TODO: add Vallet2023 to references.bib
\end{longtblr}
\vspace{2pt}
\begin{flushleft}
\footnotesize \textit{*Patients refers to the number of patients with 16S microbiome
measurements included here. Main results give a brief description of the primary
findings of the original study.}
\end{flushleft}


% =====================================================================
% TABLE 2 — Sample generation / sequencing characteristics
% Columns: Reference | Platform | Region | Primers | Reads | Avg. depth
% (Platform = Illumina/454 pyro/etc.; Region = V4, V3-V4, etc.;
%  Reads = paired / single-end; Avg. depth = mean reads per sample.)
% Data cells left as \tbd — fill in paper-by-paper.
% =====================================================================
\begin{table}[h]
\centering
\caption{Sample generation and 16S sequencing characteristics for each cohort.}
\begin{tblr}{
  width = \linewidth,
  colspec = {Q[l,m,0.16\linewidth] Q[l,m,0.14\linewidth] Q[l,m,0.08\linewidth] Q[l,m,0.24\linewidth] Q[l,m,0.12\linewidth] Q[l,m,0.12\linewidth]},
  row{1} = {font=\bfseries},
  hline{1,2,Z} = {0.8pt},
  cells = {font=\small},
  column{1} = {font=\footnotesize},
}

Reference & Platform & Region & Primers & Reads & Avg. depth (Original/ Retained) \\
Artacho2024 & Illumina MiSeq & V1-V3 & 341 F; 805 R & Paired-end & (15,3674.6/ \tbd) \\
DAmico2019 & Illumina MiSeq & V3-V4 & 341 F; 785 R & Single-end pre-merged & (23,431.2/ \tbd) \\
Fujimoto2024 & Illumina MiSeq & V3-V4 & 341 F; 785 R & Paired-end & (105,943.9/ \tbd) \\
Ingham2019 & Roche 454 & V4-V5 &  519 F; 926 R & Single-end pyro & (8,141.5/ \tbd) \\
Jarosch2023 & Ion Torrent & V1-V3 & S-D-Bact-0008-c-S-20 F; S-D-Bact-0517-a-A-18 R & Single-end pyro & (61,285.5/ \tbd) \\
Liu2017 & Illumina MiSeq & V4 & 534 F; 806 R & Single-end & (288,174.7/ \tbd) \\
Spindelboeck2021 & Ion Torrent  & V4 & 515-S3 F; 806-S2 R & Single-end pyro & (91,788.4/ \tbd) \\
Vallet2023 & Illumina MiSeq & V3-V4 & 343 F; 800R & Paired-end & (52,013.6/ \tbd) \\
\end{tblr}
\vspace{2pt}
\begin{flushleft}
\footnotesize \textit{Platform: sequencing platform (e.g.\ Illumina MiSeq, 454
pyrosequencing). Region: 16S hypervariable region targeted (e.g.\ V4, V3--V4).
Reads: paired-end or single-end. Avg.\ depth: mean reads per sample after
demultiplexing (original) and after trimming and denoising (retained).}
\end{flushleft}
\end{table}

\FloatBarrier
\newpage

% =====================================================================
% Per-study sections. Each has three subsections:
%   1. Summary            — short summary of the paper
%   2. Identified species — microbial species discussed in the original paper
%   3. Processing          — pipeline steps applied, primer sequences, layout, etc.
% Fill in paper-by-paper.
% =====================================================================
\section{Artacho2024}
\subsection{Summary}
This prospective cohort study used multi-omics to characterize gut microbiome changes during allo-HSCT and examine which changes preceded moderate-to-severe acute GVHD. The authors combined 16S rRNA sequencing, shotgun metagenomics, targeted and untargeted fecal metabolomics, machine-learning feature selection, and bacterial-metabolite network analysis. Across the transplant period, patients showed reduced bacterial richness and diversity, broad depletion of commensal \taxClostridiales, loss of microbiome-encoded functions involved in short-chain fatty acid production, reduced fecal acetate, propionate, and butyrate, and expansion of several \taxStaphylococcusSpecies. Patients who subsequently developed moderate-to-severe acute GVHD had more pronounced depletion of specific \taxClostridiales taxa, greater expansion of \taxEnterococcusFaecium and \taxParabacteroides, greater reductions in acetate and malonate, more extensive loss of microbiome-encoded functions, and more disrupted relationships between commensal bacteria and their metabolic products. In contrast, the overall decline in $\alpha$-diversity was similar between future GVHD and non-GVHD groups, indicating that changes in particular taxa, functions, metabolites, and cross-omic interactions were more informative than the global diversity change alone.
\par

The full 16S cohort comprised 172 samples across 105 allo-HSCT recipients. The first stool sample was collected, on average, approximately 4-5 days before stem-cell infusion, and the second was collected around day~+14, near engraftment. All analyzed post-transplant samples were collected before GVHD developed. The principal paired 16S analyses used 67 recipients with matched pre-transplant and engraftment samples, while paired shotgun metagenomics and metabolomic analyses were performed in smaller subsets. For the main GVHD comparison, patients who later developed grade II--IV acute GVHD were compared with patients who developed grade I or no GVHD. The authors calculated within-patient changes between the two time points and used these dynamics to identify taxonomic, functional, and metabolomic features associated with future GVHD; they also examined associations with GVHD grade, lower-gastrointestinal involvement, skin-only GVHD, and bacterial-metabolite interaction networks.
\par

\subsection{Identified microbial taxa}
Most GVHD-associated findings below refer to the change from the pre-transplant sample collected shortly before stem-cell infusion to the sample collected around engraftment, rather than abundance at either time point alone. Both samples preceded clinical GVHD.

\begin{itemize}

\item \taxClostridiales: Members of the \taxClostridiales order were broadly depleted after allo-HSCT, including prominently discussed genera such as \taxBlautia, \taxFaecalibacterium, and \taxOscillibacter. The depletion was more pronounced in patients who subsequently developed grade II--IV acute GVHD.
  \begin{itemize}
  \item \taxRuminococcaceae, \taxClostridiumIV, and \taxFlavonifractor were among the 16S-derived taxa that declined more in patients who later developed GVHD.
  \item Shotgun sequencing similarly identified greater depletion of \taxAnaeromassilibacillusSp, \taxButyricicoccusSp, and an \taxLachnoanaerobaculumSpecies in patients who later developed GVHD.
  \end{itemize}

\item \taxEnterococcusFaecium: \taxEfaecium expanded to a greater extent between the pre-transplant and engraftment samples in patients who subsequently developed moderate-to-severe acute GVHD.

\item \taxParabacteroides: \taxParabacteroides increased after allo-HSCT in patients who later developed GVHD but decreased in the comparison group. Greater expansion was also associated with more severe GVHD.

\item \taxStaphylococcusAureus, \taxStaphylococcusEpidermidis, and \taxStaphylococcusSimulans: These species expanded after allo-HSCT across the cohort and were not presented as specifically associated with future GVHD. Post-transplant \taxStaphylococcus abundance was inversely associated with fecal short-chain fatty acid levels, and the lower short-chain fatty acid concentrations measured after transplant were less inhibitory to \taxSaureus and \taxSepidermidis growth in vitro than the concentrations measured before transplant.

\item \taxFaecalibacterium, \taxClostridiumXIVb, and butyrate: \taxFaecalibacterium abundance was positively associated with butyrate before transplantation in both outcome groups. This relationship was lost after transplantation in patients who did not develop moderate-to-severe GVHD and reversed to a negative association in patients who did. Within the post-transplant GVHD interaction network, \taxClostridiumXIVb was negatively associated with \taxFaecalibacterium and positively associated with butyrate. These were bacterial-metabolite network findings rather than evidence that \taxClostridiumXIVb itself expanded in patients who developed GVHD.

\item GVHD severity and anatomical site:
  \begin{itemize}
  \item Greater depletion of \taxButyricicoccusBare showed a borderline association with more severe GVHD.
  \item Selected \taxBacteroides ASVs were positively associated with GVHD severity, but shotgun sequencing did not confirm a specific \taxBacteroidesSpecies as positively associated with severity.
  \item The pattern of greater \taxClostridiales depletion and greater \taxEfaecium and \taxParabacteroides expansion was most evident among patients who developed lower-gastrointestinal GVHD.
  \item \taxVeillonellaceae and its corresponding order \taxSelenomonadales showed distinct dynamics in the limited analysis of patients who developed skin-only GVHD.
  \end{itemize}

\end{itemize}

\subsection{Processing}
\begin{itemize}
\item \textbf{Fetch:} SRA reads were fetched from BioProject \texttt{PRJNA1088414} (104 SRR runs).
\item \textbf{Import:} Reads were imported into QIIME 2 as demultiplexed paired-end sequences.
\item \textbf{Trim:} QIIME 2 cutadapt was used to remove the amplification primers. The paper cites the Illumina ``16S Metagenomic Sequencing Library Prep'' protocol rather than printing primer sequences directly; the primers below are the Klindworth (2013) 341F/805R pair used by that protocol, and their lengths (17/21 nt) match the paper's reported positional trims.
  \begin{itemize}
  \item Forward primer (341F): \texttt{CCTACGGGNGGCWGCAG}
  \item Reverse primer (805R): \texttt{GACTACHVGGGTATCTAATCC}
  \end{itemize}
\item \textbf{Denoise:} QIIME 2 DADA2 \texttt{denoise-paired} was used.
  \begin{itemize}
  \item \texttt{--p-trim-left-f}: 0
  \item \texttt{--p-trim-left-r}: 0
  \item \texttt{--p-trunc-len-f}: 265
  \item \texttt{--p-trunc-len-r}: 225
  \end{itemize}
\end{itemize}


\section{DAmico2019}
\subsection{Summary}

This longitudinal study examined whether enteral nutrition (EN), compared with parenteral nutrition (PN), supported structural and functional recovery of the gut microbiome after pediatric allo-HSCT. Patients in the PN group were chosen based on matching to the EN group on age, source of stem cell, type of disease, and conditioning regimen. Both groups experienced an early loss of microbiota diversity and substantial compositional disruption following transplantation. However, patients receiving EN subsequently recovered microbiota diversity and a community structure resembling their pre-HSCT profiles, whereas patients receiving PN maintained lower diversity and continued to diverge from their pre-HSCT composition through follow-up. Fecal acetate, propionate, and butyrate also returned toward pre-HSCT levels in the EN group but remained depleted in the PN group. Documented bloodstream infections occurred only in the PN group. Acute GVHD occurred in the same proportion of patients in both nutritional groups, so this study did not demonstrate that EN reduced aGVHD incidence.
\par

The whole cohort of all 20 pediatric allo-HSCT recipients was used for the microbiome primary analysis. Inclusion required a stool sample collected before HSCT and at least two samples collected afterward. In total, 104 longitudinal stool samples underwent 16S rRNA sequencing and 99 underwent short-chain fatty acid analysis. The manuscript labels the initial sample as ``pre-HSCT'' but does not define a uniform collection window relative to conditioning or stem-cell infusion. The first post-transplant sample was collected between days~+8 and~+30, and later samples were collected through day~+120. The authors compared diversity, community composition, individual taxon trajectories, and fecal short-chain fatty acid production between the EN and PN groups across these sampling periods.
\par

\subsection{Identified microbial taxa}

The published taxonomic analyses compared microbiome recovery between the enteral- and parenteral-nutrition groups. They did not directly test individual taxa for association with aGVHD development. The findings below therefore describe nutrition-associated post-HSCT trajectories.

\begin{itemize}

\item \taxLachnospiraceae, \taxBlautia, and \taxDorea: \taxLachnospiraceae remained relatively stable during post-HSCT recovery in patients receiving EN but decreased in patients receiving PN. Post-HSCT differences between the nutritional groups were attributed in part to greater representation of \taxBlautia and \taxDorea in EN recipients.

\item \taxParabacteroides and \taxOscillospira: Both genera were more represented after HSCT in patients receiving EN than in patients receiving PN.

\item \taxBacteroidaceae and \taxBacteroides: \taxBacteroidaceae was more represented after HSCT in EN recipients than in PN recipients, largely because of differences in \taxBacteroides.

\item \taxFaecalibacterium: Relative abundance decreased from pre-HSCT levels during post-transplant follow-up only in both sets of patients, but the decrease was only significant in the PN group.

\item \taxStreptococcus: Relative abundance increased after transplantation in patients receiving PN but did not significantly change in patients receiving EN.


\end{itemize}

\subsection{Processing}
\begin{itemize}
\item \textbf{Fetch:} SRA reads were fetched from BioProject \texttt{PRJNA592853}.
  \begin{itemize}
  \item Reads are deposited pre-joined: each run is a single, near full-length V3-V4 amplicon (341F at the 5' end, revcomp(785R) at the 3' end) rather than separate R1/R2 pairs. Splitting the reads back into paired R1/R2 for DADA2 was attempted, but the recovered reverse reads were too short to reliably overlap on the amplicon, so the reads were instead kept single-end/pre-merged and denoised with Deblur.
  \end{itemize}
\item \textbf{Import:} Reads were imported into QIIME 2 as demultiplexed single-end sequences.
\item \textbf{Trim:} QIIME 2 cutadapt was used to remove the amplification primers.
  \begin{itemize}
  \item Forward primer (341F): \texttt{CCTACGGGNGGCWGCAG}
  \item Reverse primer, revcomp (785R): \texttt{GGATTAGATACCCBDGTAGTC}
  \end{itemize}
\item \textbf{Denoise:} QIIME 2 Deblur \texttt{denoise-16S} was used.
  \begin{itemize}
  \item \texttt{--p-trim-length}: 400
  \item \texttt{--p-left-trim-len}: 0
  \item Quality-filter minimum q-score: 4
  \end{itemize}
\end{itemize}

\section{Fujimoto2024}
\subsection{Summary}
This prospective translational study investigated why intestinal \taxEnterococcusFaecalis persists during allo-HCT and whether an \taxEfaecalis-specific phage-derived endolysin could reduce \taxEfaecalis-associated acute GVHD. Longitudinal microbiome profiling showed frequent \taxEnterococcus domination in patients, and cytolysin-poitive  \taxEfaecalis was associated with the development of aGVHD. Whole-genome and prophage analysis of the patient isolates identified an \taxEfaecalis-specific endolysin that lysed patient-derived \taxEfaecalis and its biofilms but did not lyse \taxEfaecium or other tested intestinal isolates. In gnotobiotic GVHD mouse models, the endolysin reduced intestinal \taxEfaecalis or cytolysin abundance and improved survival without consistently altering the broader bacterial community. The endolysin was not administered to patients.

The human 16S cohort comprised 46 transplant episodes from 45 adults because one patient was enrolled during two separate allo-HCT procedures. Each included transplant episode contributed at least two sequential fecal samples, for a total of 317 samples, of which 315 were viable. The pre-transplant sample was collected within 14 days before conditioning began. Subsequent samples were scheduled weekly around stem-cell infusion, defined as day 0, through approximately day 98 or hospital discharge. The authors used 16S sequencing to describe longitudinal community composition and defined \taxEnterococcus domination as greater than 25\% genus-level relative abundance, then modeled time to the first dominated sample. They subsequently cultured dominated samples to isolate \taxEfaecalis and \taxEfaecium, assessed antimicrobial susceptibility, cytolysin genes and haemolysis, and performed isolate whole-genome and prophage analyses. Selected isolates and patient-derived fecal communities were then used for the in vitro and gnotobiotic mouse experiments.

\subsection{Identified microbial taxa}
The 16S analysis established genus-level \taxEnterococcus domination. Species identification, cytolysin status, antimicrobial susceptibility, and biofilm-related functional findings came from cultured isolates, PCR assays, and isolate whole-genome sequencing rather than from 16S classification alone.

\begin{itemize}
    \item  \taxEnterococcus: Study-defined intestinal domination developed commonly during the longitudinal allo-HCT sampling period. 
    \item \taxEnterococcusFaecalis: Patient isolates carried cytolysin-associated genes, exhibited functional haemolytic activity, and showed greater representation of biofilm-associated functions than \taxEfaecium, despite remaining susceptible to the tested anti-enterococcal agents. Among transplant cases with \taxEnterococcus-dominated microbiota, fecal cytolysin positivity was associated with aGVHD. 
    \item \taxEnterococcusFaecium: Patient isolates lacked the cytolysin-associated genes detected in \taxEfaecalis and were also susceptible to the tested anti-enterococcal agents. 
\end{itemize}

\subsection{Processing}
\begin{itemize}
\item \textbf{Fetch:} SRA reads were fetched from BioProject \texttt{PRJNA1109929}
\item \textbf{Import:} Reads were imported into QIIME 2 ad demultiplexed paired-end sequences.
\item \textbf{Trim:} QIIME 2 cutadapt was used to remove the amplification primers. The primers are preceded by the TruSeq illumina overhang left over from the second-stage indexing PCR.
  \begin{itemize}
  \item Forward primer (341 F): ACACGACGCTCTTCCGATCT CCTACGGGNGGCWGCAG
  \item Reverse primer (785 R): GACGTGTGCTCTTCCGATCT GACTACHVGGGTATCTAATCC
  \end{itemize}
\item \textbf{Denoise:} QIIME 2 DADA \texttt{denois-paired} was used.
  \begin{itemize}
    \item \texttt{--p-trim-left-f}: 0
  \item \texttt{--p-trim-left-r}: 0
  \item \texttt{--p-trunc-len-f}: 265
  \item \texttt{--p-trunc-len-r}: 194
  \end{itemize}
\end{itemize}

\section{Ingham2019}
\subsection{Summary}

This longitudinal study followed children undergoing myeloablative conditioning and allo-HSCT. At each time point, information was recorded on gut microbiome, immune markers, peripheral immune-cell reconstitution, antibiotic exposure, acute GVHD severity, and survival. Microbial $\alpha$-diversity declined after transplantation, accompanied by reductions in \taxRuminococcaceae and \taxLachnospiraceae and an increase in \taxEnterococcaceae. Integrated multivariate analyses identified three major patterns. Greater \taxLactobacillaceae abundance, represented primarily by \taxLactobacillusSpp, clustered with elevated human beta-defensin 2 and monocyte counts before HSCT, moderate-to-severe aGVHD, and mortality. Greater \taxRuminococcaceae and \taxLachnospiraceae abundance clustered with more rapid NK- and B-cell reconstitution, no or mild aGVHD, and survival. \taxEnterobacteriaceae and \taxStaphylococcaceae clustered with elevated C-reactive protein and systemic inflammation but were not associated with aGVHD grade. The study did not establish that the observed microbial patterns caused or consistently preceded aGVHD.
\par

The full immune-monitoring cohort comprised 37 pediatric recipients, of whom 30 contributed 97 longitudinal fecal samples for 16S rRNA analysis. Samples were collected at up to seven time points: a ``pre-HSCT'' sample collected between day~$-33$ and day~$-3$, a sample around graft infusion between day~$-2$ and day~$+2$, and weekly samples through post-transplant week~$+5$. Conditioning began on day~$-7$, and the pre-HSCT specimens from five patients were collected after conditioning had started, so ``pre-HSCT'' did not uniformly represent a pre-conditioning baseline. Not every patient contributed a specimen at every planned time point. The authors related longitudinal OTU abundances to immune markers, immune-cell counts, antibiotic exposure, aGVHD severity, and survival using sparse partial least-squares regression and canonical correspondence analysis. They additionally assigned samples to community state types to examine whether particular community structures persisted or changed over time.
\par

\subsection{Identified microbial taxa}

The principal taxonomic results arose from integrated multivariate analyses and community-state clustering rather than from independent single-taxon association tests. These findings are exploratory associations across repeated samples. In particular, temporal analyses showed that some microbial changes occurred after aGVHD onset and therefore should not automatically be interpreted as predictors of aGVHD.

\begin{itemize}

\item \taxLactobacillaceae and \taxLactobacillusSpp: Greater abundance clustered with elevated human beta-defensin 2 and monocyte counts before HSCT, moderate-to-severe aGVHD, and mortality. A community state dominated by \taxLactobacillaceae was observed more frequently in patients with moderate-to-severe aGVHD than in those with no or mild aGVHD.
  \begin{itemize}
  \item Temporal patterns were heterogeneous. In non-survivors with moderate-to-severe aGVHD, the major expansion often occurred after aGVHD onset, whereas some survivors with moderate-to-severe aGVHD already had high abundance before aGVHD onset. The study therefore did not establish a uniform temporal sequence between this family and aGVHD.
  \end{itemize}

\item \taxRuminococcaceae, \taxLachnospiraceae, \taxFaecalibacterium, and \taxBlautia: OTUs from these families and genera were associated with more rapid NK- and B-cell reconstitution, no or mild aGVHD, and survival. A community state dominated by \taxRuminococcaceae and \taxLachnospiraceae persisted after transplantation only among patients with no or mild aGVHD.

\item \taxEnterococcaceae and \taxEnterococcusSpp: \taxEnterococcaceae increased after HSCT and was the most abundant family across the study population. However, an \taxEnterococcaceae-dominated community state was not restricted to adverse outcomes and often persisted in patients with no or mild aGVHD. Two \taxEnterococcusSpp OTUs also clustered with higher NK- and B-cell counts, indicating that the findings did not support treating all enterococcal abundance as uniformly adverse in this cohort.

\item \taxEnterobacteriaceae and \taxStaphylococcaceae: OTUs from these facultative anaerobic families clustered with elevated post-HSCT C-reactive protein and systemic inflammation. This inflammatory microbial pattern was not associated with aGVHD grade.

\item Antibiotic-associated microbial pattern: Vancomycin and ciprofloxacin exposure coincided with lower \taxRuminococcaceae and \taxLachnospiraceae abundance and higher \taxEnterobacteriaceae abundance. Antibiotic exposure therefore represented an important treatment-related influence on the observed microbial associations.

\end{itemize}

\subsection{Processing}
\begin{itemize}
\item \textbf{Fetch:} ENA reads were fetched from BioProject \texttt{PRJEB25221}. This study is 454 pyrosequencing data with several ENA/\texttt{fasterq-dump} quirks (a byte-identical duplicate file per run, mixed forward/reverse orientation tagged in the read name, and a retained 3' reverse primer plus technical tail), so a mandatory local cleanup step runs on the raw fetch output before import.
  \begin{itemize}
  \item \texttt{chpc/lib/prep\_ingham\_454.sh} normalizes each run: it drops the duplicate \texttt{\_2} file, pools the \texttt{\_1} and bare FASTQ, trims the retained 3' reverse-primer sequence and 454 technical tail with cutadapt (offering revcomp(926R) and revcomp(519F) as alternative 3' adapters so each read is trimmed in point native orientation, keeping untrimmed reads), reverse-complements reads tagged \texttt{.R2\_} in the read name so all reads end up forward-oriented (519F$\rightarrow$926R), deduplicates by read ID, and writes one clean, primer-free, single-orientation FASTQ per run (archiving the raw inputs to \texttt{RawData/\_raw454/}).
  \end{itemize}
\item \textbf{Import:} Reads were imported into QIIME 2 as demultiplexed single-end sequences.
\item \textbf{Trim:} Primers were already removed by \texttt{prep\_ingham\_454.sh} at the fetch step above, so the standard QIIME 2 cutadapt 'trim' stage is a no-op for this study.
  \begin{itemize}
  \item Forward primer (519F): \texttt{CAGCAGCCGCGGTAATAC}
  \item Reverse primer (926R): \texttt{CCGTCAATTCCTTTGAGTTT}
  \end{itemize}
\item \textbf{Denoise:} QIIME 2 DADA2 \texttt{denoise-pyro} was used.
  \begin{itemize}
  \item \texttt{--p-trim-left}: 0
  \item \texttt{--p-trunc-len}: 368
  \item \texttt{--p-max-ee}: 2.0
  \item \texttt{--p-trunc-q}: 2
  \end{itemize}
\end{itemize}

\section{Jarosch2023}
\subsection{Summary}

This post-transplant multimodal observational study investigated how antibiotic-associated gut microbiome disruption relates to intestinal regulatory T-cell differentiation and the severity of GI-GVHD after allo-HSCT. The authors combined biopsy-matched stool 16S profiles with single-cell RNA and paired T-cell receptor sequencing of gastrointestinal biopsies, multiplexed tissue imaging, clinical biomarkers, and treatment data. Severe GI-GVHD was associated with dense infiltration and clonal expansion of activated CD8 T cells, and expanded clones persisted over time and were detected across anatomically distant gastrointestinal sites. Broad-spectrum antibiotic exposure was associated with lower microbiome diversity, reduced intestinal Treg infiltration, and Tregs with less differentiated and less suppressive phenotypes. SCFA-producing taxa, particularly Clostridia (class) and Bacteroidia (class), were associated with more suppressive Treg phenotypes, whereas Gammaproteobacteria (class) were enriched among taxa associated with higher T-effector gene scores in Tregs. Overall 16S $\beta$-diversity did not cluster according to GI-GVHD severity, antibiotic exposure, or time after transplantation, indicating that the microbial findings reflected associations between particular taxa and immune phenotypes rather than a broad community-level separation. 
\par

The full multimodal tissue cohort included 60 allo-HSCT recipients. The 16S component comprised 52 stool samples collected within the week of a corresponding gastrointestinal biopsy, most commonly on the same day; the main text did not report a separate unique-recipient count for these microbiome samples. The single-cell RNA and T-cell receptor dataset included 31 biopsies from 26 recipients, while the ChipCytometry dataset included 52 biopsies from 40 recipients, with overlap between the analytical cohorts. Biopsies were collected during screening in recipients without clinical signs of GVHD, within seven days of the first GVHD symptoms, or more than seven days after symptom onset. The microbiome samples therefore do not represent a common baseline and may reflect pre-onset, onset, or established disease. For the taxon-level analysis, recipients were grouped according to the presence or absence of each taxon, and microbial features were related to Treg, T-effector, and Treg-suppression gene scores.
\par

\subsection{Identified microbial taxa}
\tbd

\subsection{Processing}
\begin{itemize}
\item \textbf{Fetch:} ENA reads were fetched from BioProject \texttt{PRJEB60178} (53 single-end Ion Torrent runs).
\item \textbf{Import:} Reads were imported into QIIME 2 as demultiplexed single-end sequences.
\item \textbf{Trim:} QIIME 2 cutadapt was used to remove the amplification primers. The reverse primer sequence below is reconstructed from the Klindworth (2013) probe name rather than taken directly from the paper and still needs confirmation against the deposited reads.
  \begin{itemize}
  \item Forward primer (27F / S-D-Bact-0008-c-S-20): \texttt{AGRGTTYGATYMTGGCTCAG}
  \item Reverse primer, revcomp (S-D-Bact-0517-a-A-18): \texttt{GCCAGCMGCCGCGGTAAT}
  \end{itemize}
\item \textbf{Denoise:} QIIME 2 DADA2 \texttt{denoise-pyro} was used.
  \begin{itemize}
  \item \texttt{--p-trim-left}: 15 (per the authors' reported 5' trim)
  \item \texttt{--p-trunc-len}: \textcolor{red}{FILL}
  \item \texttt{--p-max-ee}: 5.0 (per the paper)
  \item \texttt{--p-trunc-q}: 2
  \end{itemize}
\end{itemize}


\section{Liu2017}
\subsection{Summary}

This prospective single-center study asked whether gut microbiota features present before conditioning in allo-HSCT recipients, microbiota dissimilarity between recipients and their HLA-matched sibling donors, or the donors' own microbiota were associated with subsequent acute gastrointestinal GVHD (GI-GVHD). The authors also evaluated baseline recipient microbiota in relation to overall mortality. Compared with healthy sibling donors, recipients had lower phylogenetic diversity, different community membership, and frequent disturbance-associated domination by facultative anaerobes compared to obligate anaerobes. Greater recipient baseline diversity was associated with lower overall mortality, but recipient diversity and donor-recipient dissimilarity were not significantly associated with GI-GVHD. In the paired-donor subcohort, greater donor diversity was associated with lower odds of GI-GVHD in the recipient. No individual recipient or donor taxon remained significantly associated with GI-GVHD after multiple-testing correction, making the taxa-level findings exploratory.
\par

The cohort comprised 57 allo-HSCT recipients, 22 of whom had a paired HLA-matched sibling donor sample. Recipient ``baseline'' was defined as a fecal sample collected 0--7 days before preparative conditioning; donor baseline was collected 0--7 days before hematopoietic growth-factor administration. The full recipient cohort was used to compare recipient and donor community features, to model recipient baseline diversity and species-level relative abundances against GI-GVHD, and to test baseline diversity against overall mortality. The 22 paired donor-recipient samples were additionally used to test whether donor-recipient UniFrac dissimilarity, donor diversity, or donor taxa predicted GI-GVHD.
\par


\subsection{Identified microbial taxa}

The manuscript reported 58 species-level differences between recipients and donors in a supplementary table. The list below is restricted to taxa explicitly emphasized by the authors in the main text as differing between recipients who did or did not develop GI-GHVD, although the associations, which were found by multiple logistic regression, did not persist after false discovery rate correction. 
\begin{itemize}
\item \taxParabacteroidesDistasonis and \taxBarnesiellaceaeSpp: Higher baseline (pre-conditioning) relative abundance was associated with lower odds of subsequent GI-GVHD.
\item \taxLachnobacteriumSpp and \taxHaemophilusParainfluenzae: Higher baseline (pre-conditioning) relative abundance was associated with higher odds of subsequent GI-GVHD 

\end{itemize}


\subsection{Processing}
\begin{itemize}
\item \textbf{Fetch:} ENA reads were fetched from European Nucleotide Archive ERP017899
\begin{itemize}
    \item R1 and R2 reads were uploaded separately. Preprocessing was done to determine which read was which orientation. As only the V4 region was sampled, foreward reads carried the full length of the sample, and only forward reads were kept. 
\end{itemize}
\item \textbf{Import:} Reads were imported into QIIME 2 as demultiplexed single-end sequences. 
\item \textbf{Trim:} QIIME 2 cutadapt was used to remove the amplification primers
  \begin{itemize}
  \item Forward primer (534F/515F Fwd): GTGCCAGCMGCCGCGGTAA
  \item Reverse primer (806R Rev): GGACTACHVGGGTWTCTAAT
  \end{itemize}
\item \textbf{Denoise:} QIIME 2 DADA2 \texttt{denoise-single} was used.  
  \begin{itemize}
  \item \texttt{--p-trunc-len}: 230
  \item \texttt{--p-max-ee}: 2
  \end{itemize}
\end{itemize}


\section{Spindelboeck2021}
\subsection{Summary}

This prospective single-center study evaluated repeated colonoscopically delivered fecal microbiota transplantation (FMT) as salvage therapy for severe, treatment-refractory lower gastrointestinal acute graft-versus-host disease (GI-GVHD) after allo-SCT. The authors sought to determine clinical response, identify factors associated with response, and relate microbiota restoration to changes in ileocolonic mucosal immune cells. The primary endpoint was sustained remission of GI-GVHD 90 days after FMT. Of the 10 patients, four reached the endpoint and were labeled 'responders', five did not reach the endpoint and were labeled 'non-responders', and one patient died before the 90 days. Responders had significantly longer overall and GI-GVHD-related survival. Response occurred only among patients in whom broad-spectrum antibiotics could be withheld for substantial periods. Responders developed donor-like $\alpha$- and $\beta$-diversity, stronger genus-level similarity to donors, and markedly greater predicted donor-microbiota engraftment, whereas nonresponders retained low-diversity dysbiosis. Responders were enriched for anaerobic, predominantly butyrate-producing taxa. At baseline all patients had similar mucosal immune profiles, but post-FMT responders shifted toward fewer CD8+ and Th17 cells and more regulatory T cells and type 3 innate lymphoid cells.
\par

The human 16S cohort comprised all 10 adults with severe treatment-refractory GI-GVHD who received FMT as well as 14 donor samples that were taken prior to donation. For patients, stool microbiota were profiled before the first FMT (baseline), during the weekly series of up to six FMTs, and after treatment. For cross-sectional $\alpha$- and $\beta$-diversity analyses, the authors used 10 baseline patient samples, the last available post-last-FMT sample from each patient (four responders and six nonresponders/other), and 14 donor samples. The authors compared observed-species $\alpha$-diversity, weighted UniFrac $\beta$-diversity, genus-level correlations, LEfSe-discriminatory taxa, SourceTracker-estimated donor contribution, and samples collected with versus without antibiotic exposure during the preceding seven days to test whether clinical response tracked with donor microbiota engraftment and whether antibiotic use impeded microbiome recovery.
\par

\subsection{Identified microbial taxa}
Differences in microbial taxa were analyzed by relative abundance pre-FMT in patients compared to health donors, differential abundance in final post-FMT samples between resonders and non-responders, and identification as a discriminator of non-response by LEfSe. 
\begin{itemize}
\item \taxEnterococcus: Overabundant in patients before FMT relative to healthy donors. Post-FMT, the two patients had enterococcus domination were in the non-responder group.
\item \taxLactobacillus: Overabundant in patients before FMT relative to healthy donors. Significantly more abundant in nonresponders than compared to responders. LEfSe identified it as a discriminator of nonresponse.
\item \taxOscillibacter, \taxSubdoligranulum, and \taxButyricicoccus (belonging to family \taxRuminococcaceae), \taxBlautia, \taxLachnoclostridium, \taxCoprobacter, \taxParabacteroides, \taxClostridiumInnocuumGroup, and \taxAlistipes: Responder-associated genera in the LEfSe analysis.
\begin{itemize}
    \item The authors describe these responder enriched genera as anaerobes, and point out that \taxBlautia and \taxRuminococcaceae are butyrate-producing. 
\end{itemize}
\item \taxCampylobacter and \taxStaphylococcus: overabundant in patients before FMT relative to healthy donors in the baseline LEfSe analysis.
\end{itemize}
\subsection{Processing}
\begin{itemize}

\item \textbf{Fetch:} SRA reads were fetched from BioProject \texttt{PRJEB39834}.
\item \textbf{Import:} Reads were imported into QIIME 2 as demultiplexed single-end sequences.
\item \textbf{Trim:} QIIME 2 cutadapt was used to remove the amplification primers.
  \begin{itemize}
  \item Forward primer (515-S3 Fwd): \texttt{TGCCAGCAGCCGCGGTAA}
  \item Reverse primer (806-S2 Rev): \texttt{GGACTACCAGGGTATCTAAT}
  \end{itemize}
\item \textbf{Denoise:} QIIME 2 DADA2 \texttt{denoise-pyro} was used.
  \begin{itemize}
  \item \texttt{--p-trunc-len}: 257
  \item \texttt{--p-trim-left}: 0
  \item \texttt{--p-max-ee}: 2
  \item \texttt{--p-trunc-q}: 2
  \end{itemize}
\end{itemize}

\section{Vallet2023}
\subsection{Summary}
This longitudinal multi-omic substudy of the multicenter, randomized, double-blind, placebo-controlled ALLOZITHRO trial (\cite{Bergeron2017}) examined whether early azithromycin exposure was associated with changes in the post-allo-HSCT gut bacteriome, virome, and metabolome and whether those changes were associated with hematologic relapse. The ALLOZITHRO trial itself was canceled early due to safety reasons, with the arm receiving azithromycin experiencing a higher risk or relapse and death. In this substudy, the authors identified four bacterial enterotypes with distinct diversity and taxonomic profiles. Enterotype 1 was overrepresented among samples from the azithromycin arm, whereas enterotype 2 was overrepresented among placebo samples and was strongly associated with complete remission at 12 months. Enterotypes 1, 3, and 4 contained larger proportions of samples from patients who relapsed. Patients whose enterotype changed during the early transplant period were more often in complete remission at 12 months than patients whose enterotype remained stable. A phylotype related to Bacteroides fragilis was associated with greater relapse risk, whereas phylotypes related to Bacteroides sp. DJF\_B097 and Prevotella sp. DJF\_RP53 were associated with lower relapse risk. Enterotypes were also associated with fecal and plasma metabolite profiles and circulating T-cell states, including an association between the B. fragilis-related phylotype and exhausted T-cell phenotypes. The gut virome as a whole was not significantly associated with treatment arm, relapse, or bacterial enterotype, although individual bacteriophages correlated with specific bacterial phylotypes. Enterotype distribution was not associated with acute or chronic GVHD.

The human microbiome cohort comprised 55 allo-HSCT recipients, including 27 assigned to azithromycin and 28 assigned to placebo. The study included 148 longitudinal fecal samples—73 from the azithromycin arm and 75 from the placebo arm—with the sampling window beginning during the week before allo-HSCT and continuing through six weeks after transplantation. Of these, 16S sequencing data was available for 118 samples. The authors clustered sample-level 16S profiles into enterotypes, tested their associations with treatment arm and relapse or complete remission at 12 months, and used PERMANOVA and longitudinal microbial-signature analyses to identify contributing phylotypes. Relapse risk for individual phylotypes was subsequently tested with Fine-Gray competing-risk models using the mean OTU abundance for each patient. Cross-omic correlation networks connected leading bacterial phylotypes with fecal metabolites and viral species, and selected enterotypes and outcome-associated phylotypes were compared with plasma metabolomic and circulating T-cell-state data.

\subsection{Identified microbial taxa}
The main results of the paper were presented by bacterial Enterotypes.
\begin{itemize}
    \item Enterotype 1: Characterized primarily by \textit{Clostridium XIVa} and \textit{Parabacteroides}. Had the second-highest alpha diversity of the four clusters. 32\% of samples with this enterotype relapsed. Was not associated with greater remission after 12 months.
    \item Enterotype 2: Characterized by \textit{Bacteroides}, \textit{Fusobacterium}, and \textit{Faecalibacterium}. Had the highest alpha diversity. Compared with the other Bacteroides-rich enterotype, enterotype 3, it was enriched in phylotypes related to \textit{Bacteroides uniformis}, \textit{Bacteroides sp. Smarlab BioMol-2301151}, \textit{Bacteroides ovatus}, and \textit{Bacteroides massiliensis}. Associated with greater remission after 12 months. 
    \item Enterotype 3: Also enriched in \textit{Bacteroides}, but it had significantly lower richness and diversity than enterotype 2. Relative to enterotype 2, enterotype 3 was enriched in phylotypes related to \textit{Bacteroides vulgatus} and \textit{Bacteroides fragilis}. Was not associated with greater remission after 12 months.
    \item Enterotype 4: Driven primarily by \textit{Enterococcus} and \textit{Enterobacter} and had the lowest alpha diversity. Was not associated with greater remission after 12 months.
\end{itemize}

There were additional results relating to specific phylotypes.
\begin{itemize}
    \item Phylotype related to \textit{Bacteroides fragilis}: Associated with greater relapse risk. 
    \begin{itemize}
        \item Cross-omic module was associated with branched-chain amino-acid metabolism, one Larmunavirus, and two Ceduovirus species. 
        \item Its Higher abundance was also associated with exhausted central-memory CD4 Th1, CD4 Th2, and CD8 T-cell states co-expressing TIGIT, PD-1, and TOX.
    \end{itemize}
    \item Phylotype related to \textit{Prevotella sp. DJF\_RP53}: Associated with lower relapse risk.
    \begin{itemize}
        \item Cross-omic module was associated with pentose metabolism, one Sandinevirus, and one Pepyhexavirus.
    \end{itemize} 
    \item Phylotype related to \textit{Bacteroides sp. DJF\_B097}: Associated with lower relapse risk.
    \begin{itemize}
        \item Cross-omic module was associated with phospholipid and lysophospholipid metabolism and one unclassified Siphoviridae species.
        \item  Higher abundance was associated with KLRG1+2B4+activated effector-memory CD4 Th0 cells and TIGIT+central-memory CD4 Th1 cells.
    \end{itemize}
\end{itemize}

\subsection{Processing}
\begin{itemize}
\item \textbf{Fetch:} SRA reads were fetched from BioProject \texttt{PRJNA902819}
\item \textbf{Import:} Reads were imported into QIIME 2 as demultiplexed paired-end sequences.
\item \textbf{Trim:} QIIME 2 cutadapt was used to remove the amplification primers. We note here that a typo appears in the original text with a leading T mis-added to the forward primer. We print the correct reversion here. 
  \begin{itemize}
  \item Forward primer (): ACGGRAGGCAGCAG
  \item Reverse primer (): TACCAGGGTATCTAAT
  \end{itemize}
\item \textbf{Denoise:} QIIME 2 DADA2 \texttt{denoise-paired} was used.
  \begin{itemize}
  \item \texttt{--p-trim-left-f}: 0
  \item \texttt{--p-trim-left-r}: 0
  \item \texttt{--p-trunc-len-f}: 230
  \item \texttt{--p-trunc-len-r}: 225
  \end{itemize}
\end{itemize}



\newpage
\section{Bibliography}
\printbibliography[heading=bibintoc]


\end{document}
